\section[Progression-free sets in polynomial rings and finite rings]{Progression-free sets in polynomial rings and finite rings; additive constraints}\label{sec:gpp}

\providecommand{\gppA}{\mathcal{A}}
\providecommand{\gppG}{\mathcal{G}_q}
\providecommand{\gppP}{\mathcal{P}}
\providecommand{\gppPr}{\operatorname{Pr}}
\providecommand{\gppE}{\mathbb{E}}
\providecommand{\gppVar}{\operatorname{Var}}
\providecommand{\gppdens}{\operatorname{dens}}
\providecommand{\gpplc}{\operatorname{lc}}
\providecommand{\gppm}{\mathfrak{m}}
\providecommand{\gppsloppy}{\sloppy}

The source proof selects its distinguished infinite places by a finite combinatorial
construction (\cref{src:sec:mechanism}). The lemma
\code{modEq_orderOf_of_div_pow_eq} turns a multiplicative identity between powers of one
element into a weighted additive congruence between exponents, and the lemma
\code{exists_three_term_free} produces a large set of exponents satisfying no such
congruence. The second lemma is a Behrend-type construction \cite{Behrend1946Sets}: digit
vectors on a sphere, encoded in a base large enough to prevent carries. This section
collects the results obtained by moving that exponent principle to other settings. In
$\Fq[x]$ the exponent map is the whole vector of prime valuations and the congruence becomes
an equality; this gives answers to the questions and to Conjecture~6.3 of Asada, Fourakis,
Manski, McNew, Miller and Moreland \cite{Asada2017Subsets} on sets free of three-term
geometric progressions. Over finite coefficient rings the same ideas, combined with the
density Hales--Jewett theorem, answer their Question~6.5 for large families of rings. The
remaining subsections treat additive constraints coming directly from the source
construction: periodic sets of integers with a prescribed spectrum of ratios, the exact
capacity of the sparse-place condition of the source, and coefficient collisions in the
root lattice $A_n$. Nothing in this section uses \cref{src:thm:dgrh} or any other
hypothesis on zeros of $L$-functions. The corresponding questions for sets of integers and
of ideals are treated in \cref{sec:gpi}.

\subsection{The weighted digit-sphere lemma}\label{gpp:sec:digits}
{\raggedright\sources{dirichlet-grh-research; exponent_seven_and_grh_transfers; mcnew_counterexample_and_grh_transfers; grh_method_transfers_graph_towers_and_finite_rings; grh_method_transfers}}

The following is the finite estimate behind \code{exists_three_term_free}, with all digit
losses retained. It is used in
\cref{gpp:sec:rings,gpp:sec:syndetic,gpp:sec:sparse}.

\begin{lemma}[Weighted digit-sphere lemma; \stU]\label[lemma]{gpp:lem:digits}
Let $n,H,L,o$ be positive integers, put $B=L(H-1)+1$, and assume $B^n\le o$; this holds in
particular when $(LH)^n\le o$. There is a set $S\subset\set{0,\ldots,o-1}$ with
$\abs{S}\ge H^n/(nH^2)$ such that for all $u,v,w\in S$ and all positive integers $b,c$
with $b+c\le L$,
\[
 (b+c)v\equiv bu+cw \pmod o\quad\Longrightarrow\quad u=v=w .
\]
\end{lemma}

\begin{proof}
Take the digit cube $\set{0,\ldots,H-1}^n$ and its centre
$\kappa=((H-1)/2,\ldots,(H-1)/2)$. All values of $\norm{x-\kappa}^2$ differ by integers,
since $a^2-(H-1)a$ is an integer for each digit $a$, and their range has length at most
$n(H-1)^2/4$. The most populated sphere centred at $\kappa$ therefore contains at least
$H^n/(1+\floor{n(H-1)^2/4})\ge H^n/(nH^2)$ points. Encode each point $x$ of that sphere
by $e(x)=\sum_{k=0}^{n-1}x_kB^k$ and let $S$ be the set of these integers. In a weighted
relation as in the statement, each digit of either side is at most $L(H-1)=B-1$, so both
sides lie in $[0,B^n)\subseteq[0,o)$. The congruence is therefore an equality of integers
and holds digit by digit. The underlying vectors satisfy
$v=\frac{b}{b+c}u+\frac{c}{b+c}w$. On a sphere a strict convex combination of two
different vectors has smaller distance from the centre, so $u=w$ and then $v=u$.
\end{proof}

\begin{corollary}[\stU]\label[corollary]{gpp:cor:behrend}
For every integer $L\ge2$ there is a constant $C_L$ such that for each integer $m\ge2$
some $S\subset\set{0,\ldots,m-1}$ of size at least $m\exp(-C_L\sqrt{\log(2+m)})$ contains
no three elements $d_1,d_2,d_3$, not all equal, with $(a+b)d_2=bd_1+ad_3$ for positive
integers $a,b$ with $a+b\le L$.
\end{corollary}

\begin{proof}
For large $m$ apply \cref{gpp:lem:digits} with $o=m$, $n=\floor{\sqrt{\log m}}$ and
$H=\floor{m^{1/n}/L}$. Then $(LH)^n\le m$, and since $m^{1/n}/L\to\infty$ eventually
$m^{1/n}/(2L)\le H\le m^{1/n}$. Hence
\[
 \abs{S}\ge\frac{m}{n(2L)^n m^{2/n}},
\]
and the logarithm of the denominator is $\log n+n\log(2L)+2(\log m)/n=O_L(\sqrt{\log m})$.
A larger constant covers the finitely many smaller values of $m$.
\end{proof}

\subsection{Polynomial rings over finite fields: conventions}\label{gpp:sec:conventions}
{\raggedright\sources{grh-progression-densities-and-cubic-prime-race; grh-methods-polynomial-density; grh_beyond_dirichlet_2026-11-25}}

Let $q$ be a prime power. A \emph{triple} is $f,rf,r^2f$ with all three terms nonzero
polynomials in $\Fq[x]$. Two conventions for the ratio are used. Under the
\emph{polynomial-ratio} convention $r\in\Fq[x]$ and $\deg r\ge1$. Under the
\emph{rational-ratio} convention $r$ is any nonconstant element of $\Fq(x)$, including
functions of degree zero such as $x/(x+1)$. Constant ratios are ignored in both. Write
$B_N=\set{f:\deg f\le N}$, including zero, so that $\abs{B_N}=q^{N+1}$. The \emph{natural
density} of a set $E$ is the limit of $\abs{E\cap B_N}/q^{N+1}$ when it exists, written
$\gppdens(E)$, and the \emph{upper density} is the upper limit. Working with monic
polynomials and then taking all nonzero scalar multiples gives the same limiting
densities.

Let
\[
 \gppA=\set{a\ge0:\text{every ternary digit of }a\text{ is }0\text{ or }1},\qquad
 \sum_{a\in\gppA}t^a=\prod_{j\ge0}\bigl(1+t^{3^j}\bigr).
\]
The set $\gppA$ contains no nonconstant three-term arithmetic progression: in an equality
$a+c=2b$ with $a,b,c\in\gppA$ neither side has ternary carries, so at each digit the only
possibilities are $0+0=2\cdot0$ and $1+1=2\cdot1$, and $a=b=c$. The argument allows a
common difference of either sign. Put
\begin{equation}\label{gpp:eq:mq}
 m_q=(1-q^{-1})\sum_{a\in\gppA}q^{-a}=(1-q^{-2})\prod_{j\ge1}\bigl(1+q^{-3^j}\bigr).
\end{equation}
In the ball $B_N$ call $d=N-\deg f$ the \emph{deficit} of $f$. The rule $d\in\gppA$
selects a polynomial-ratio progression-free subset of $B_N$ of density tending to $m_q$,
because a nonconstant polynomial ratio gives strictly increasing degrees in arithmetic
progression. By \cite[Theorem~4.3]{Asada2017Subsets}, $m_q$ is optimal among sets defined
solely by degree. Let $\beta_q$ be the supremum of the upper densities of
polynomial-ratio progression-free subsets of $\Fq[x]$. Conjecture~6.3 of
\cite{Asada2017Subsets} asserts $\beta_q=m_q$.

Let $\gppG$ consist of the nonzero polynomials with $v_P(f)\in\gppA$ at every monic
irreducible $P$. By \cite[Theorem~2.2]{Asada2017Subsets} this is the greedy
progression-free set. Since valuations along a triple form an arithmetic progression at
each prime, with a common difference of either sign under the rational-ratio convention,
$\gppG$ contains no triple under either convention.

\begin{lemma}[Counting valuations; \stU]\label[lemma]{gpp:lem:valuations}
\gppsloppy
Set $w_P=q^{-\deg P}$ and $p_P=(1-w_P)\sum_{a\in\gppA}w_P^a$. Among monic polynomials of
exact degree $n$, the valuations at a fixed finite list of monic irreducibles have, as
$n\to\infty$, independent limiting geometric laws $\gppPr(v_P=a)=(1-w_P)w_P^a$. A set
defined by restricting the valuations at a finite set $S$ of primes to a subset of
$\Z_{\ge0}^S$ and by requiring $v_P\in\gppA$ at every $P\notin S$ has a natural density,
equal to the product of the limiting probability of the restriction at $S$ and of
$C_S=\prod_{P\notin S}p_P$. In particular $\gppdens(\gppG)=\prod_Pp_P$, and
\begin{equation}\label{gpp:eq:CS}
 C_S\ge\prod_P(1-w_P^2)=1-q^{-1}.
\end{equation}
\end{lemma}

\begin{proof}
The geometric laws follow by counting the multiples of products of prime powers and then
excluding the next powers. Unbounded valuations at a fixed finite list can be truncated
with error tending to zero. A restriction $v_P\in\gppA$ at a prime of degree greater than
$D$ can fail only if $P^2\mid f$; the total failure probability is at most
$\sum_{d>D}q^dq^{-2d}$ uniformly in $n$, using the number of monic polynomials of degree
$d$ as an upper bound for the number of primes. Thus the exact-degree limiting proportions
exist, and summing the geometrically weighted degree layers gives the same ball density.
For \eqref{gpp:eq:CS}, $0,1\in\gppA$ gives $p_P\ge(1-w_P)(1+w_P)$, and the product
follows by unique factorization from $\prod_P(1-u^{\deg P})^{-1}=(1-qu)^{-1}$ at
$u=q^{-2}$.
\end{proof}

\subsection{A counterexample to the density conjecture for polynomial ratios}\label{gpp:sec:counterexample}
{\raggedright\sources{grh-methods-polynomial-density; grh_beyond_dirichlet_2026-11-25}}

The construction changes the admissible deficits inside classes of polynomials with a
prescribed rough part. It therefore uses the factorization of polynomials of the same
degree, which no set defined by degree alone can do.

\begin{lemma}[Replacing degree deficits within rough-factor classes; \stU]\label[lemma]{gpp:lem:replacement}
Let $A_0\subseteq\Z_{\ge0}$ contain no nonconstant three-term arithmetic progression.
Suppose $D\subseteq\Z_{\ge0}$ is also free of such progressions and
$D\setminus A_0\subseteq[0,L]$ for an integer $L\ge1$. Let $\gppP$ be the set of all monic
irreducibles in $\Fq[x]$ of degree at most $L$, and let
\[
 H(u)=\prod_{P\in\gppP}(1-u^{\deg P})^{-1}=\sum_{n\ge0}H_nu^n .
\]
If, for some $0<z<1$,
\[
 J(z):=\sum_{d\in D}z^d-\sum_{a\in A_0}z^a>0,
\]
then some integer $R\ge0$ satisfies $\Delta_R=[u^R]H(u)J(u)\ge1$, and for every such $R$
there is one infinite polynomial-ratio progression-free set of upper density at least
\begin{equation}\label{gpp:eq:delta}
 (1-q^{-1})\sum_{a\in A_0}q^{-a}+(q-1)\Delta_Rq^{-R-1}c_q,
 \qquad
 c_q=(1-q^{-1})\prod_{P\in\gppP}(1+q^{-\deg P})^{-1}>0 .
\end{equation}
In particular, for $A_0=\gppA$, $\beta_q\ge m_q+(q-1)\Delta_Rq^{-R-1}c_q$.
\end{lemma}

\begin{proof}
All series converge absolutely at $z$. The product $H(u)J(u)$ has integer coefficients and
$H(z)J(z)>0$. If every coefficient were nonpositive, the value at $z$ would be
nonpositive. Thus a coefficient $\Delta_R\ge1$ exists.

Call a monic polynomial \emph{smooth} if all its irreducible factors belong to $\gppP$
and \emph{rough} if none do. Here $H_n$ counts the monic smooth polynomials of degree $n$,
by unique factorization. Every nonzero polynomial has a unique expression
\[
 f=cgh,\qquad c\in\Fq^\times,\quad g\text{ monic rough},\quad h\text{ monic smooth}.
\]
At a degree cutoff $N\ge R$, the baseline set consists of the polynomials whose deficit
$d=N-\deg f$ lies in $A_0$. It is progression-free because degrees, and hence deficits,
would otherwise form a nonconstant arithmetic progression. For each squarefree rough monic
$g$ of degree $N-R$, change the rule on the entire class $\set{cgh}$ to the deficit rule
$d\in D$, and denote the resulting finite set by $E_N$. A polynomial in this class with
deficit $d$ has $\deg h=R-d$. Consequently each $g$ changes the number of selected
polynomials by exactly $(q-1)\Delta_R$.

A new progression cannot leave a modified class. Indeed, any progression in $E_N$
contains an added polynomial $f=cgh$, whose deficit is in $D\setminus A_0$, hence at most
$L$. If $f$ is the first or the middle term, a later term is still inside $B_N$, so the
ratio has degree at most $L$ and every irreducible factor of the ratio is smooth. If $f$
is the last term, the square of the ratio divides it; its rough part $g$ is squarefree,
so again every irreducible factor of the ratio is smooth. In either case all three terms
have the same monic rough part $g$, and their three distinct deficits lie in $D$. This
contradicts the progression-free property of $D$. The argument uses the condition
$r\in\Fq[x]$ of Conjecture~6.3.

Let $G_n$ count the squarefree rough monic polynomials of degree $n$. Unique factorization
gives
\[
 \sum_{n\ge0}G_nu^n=\frac{1-qu^2}{1-qu}\prod_{P\in\gppP}(1+u^{\deg P})^{-1},
\]
by the Euler product $\prod_P(1-u^{\deg P})^{-1}=(1-qu)^{-1}$ and
$1+y=(1-y^2)/(1-y)$. Subtracting the pole part at $u=1/q$ leaves a rational function
analytic on $\abs{u}<1$, and coefficient comparison gives $G_n/q^n\to c_q$. The exact
density of $E_N$ is
\[
 \frac{\abs{E_N}}{q^{N+1}}
  =(1-q^{-1})\sum_{\substack{a\in A_0\\a\le N}}q^{-a}
      +(q-1)\Delta_R\,\frac{G_{N-R}}{q^{N+1}},
\]
which tends to \eqref{gpp:eq:delta}.

It remains to obtain a single infinite set. Choose integers $N_j$ increasing so fast that
\begin{equation}\label{gpp:eq:blocks}
 N_j-j>2N_{j-1},\qquad N_j-2j>N_{j-1},\qquad N_1\ge R+2 .
\end{equation}
From $E_{N_j}$ keep only the polynomials of degree in $[N_j-j,N_j]$, and take the union of
these blocks over $j$. A progression whose first two terms lie in earlier blocks than its
last term has last degree at most $2N_{j-1}$. A progression with two terms in the block
of its last term has first degree at least $N_j-2j$. Both possibilities contradict
\eqref{gpp:eq:blocks}, and progressions within one block were already excluded. The
deleted lower degrees have density at most $q^{-j-1}$ at the cutoff $N_j$, which tends to
zero. The union therefore has upper density at least \eqref{gpp:eq:delta}.
\end{proof}

\begin{theorem}[Counterexample to Conjecture~6.3 of \cite{Asada2017Subsets}; \stU]\label{gpp:thm:counterexample}
For every prime power $q$ there is an infinite polynomial-ratio progression-free set
$E\subseteq\Fq[x]$ whose upper density is strictly greater than $m_q$. Consequently
$\beta_q>m_q$.
\end{theorem}

\begin{proof}
Apply \cref{gpp:lem:replacement} with $A_0=\gppA$, $L=480$ and $D=B$, where
\[
\begin{split}
B=\{&
0,1,4,5,11,13,14,16,20,30,37,38,41,42,48,50,51,53,57,67,\\
&111,112,115,116,122,124,125,127,131,141,148,149,152,153,\\
&159,161,162,164,168,178,333,334,337,338,344,346,347,349,\\
&353,363,370,371,374,375,381,383,384,386,390,437,440,441,\\
&444,449,450,453,455,464,467,468,471,476,477,480\}.
\end{split}
\]
This set of $74$ deficits is the slice $878,\ldots,1358$, translated by $-878$, of the
record $a(143)\le1360$ of G.~Theobald in the database of non-averaging sets
\cite{WroblewskiNonaveraging}. Only the displayed list is used: all $\binom{74}{2}=2701$
midpoint tests were checked exactly, so $B$ contains no nonconstant three-term arithmetic
progression. No optimality of the list is used. Since $D=B$ is finite, on the modified
classes the whole original rule is replaced, and no tail of $\gppA$ has to be joined to
$B$.

Take $z=127/128$ and let
$\gppA_7=\set{\sum_{j=0}^{6}\epsilon_j3^j:\epsilon_j\in\set{0,1}}$, which has $128$
elements and maximum $1093$. Clearing the denominator $128^{1093}$ proves the integer
inequality
\[
 \sum_{b\in B}z^b-\sum_{a\in\gppA_7}z^a>\frac{16}{1000}.
\]
Every other element of $\gppA$ is at least $3^7=2187$, so
\[
 \sum_{a\in\gppA\setminus\gppA_7}z^a\le\frac{z^{2187}}{1-z}=128z^{2187}<\frac1{100000},
\]
again by integer multiplication. Hence
\begin{equation}\label{gpp:eq:Jpositive}
 J(127/128)>\frac{3}{200},
\end{equation}
and the lemma applies. \gppsloppy The finite inputs are verified by
\artifact{grh-methods-polynomial-density/check_tilted.py} and, independently, by
\artifact{grh_beyond_dirichlet_2026-11-25/code/check_density_certificate.py}.
\end{proof}

\begin{example}[A database-free witness on a digit sphere]\label[example]{gpp:ex:sphere}
There is a fully elementary alternative certificate, which uses the same carry-free
sphere as \code{exists_three_term_free}. Put $M=3^{20}$. For
$u\in\set{0,1,\ldots,5}^9$ define
\[
 b(u)=\sum_{i=0}^{8}u_i11^i,\qquad
 B_0=\Bigl\{b(u):\sum_{i=0}^{8}(2u_i-5)^2=105\Bigr\}.
\]
Each squared coordinate is $1$, $9$ or $25$, with multiplicity two. If the three counts
are $i,j,k$, then $i+j+k=9$ and $j+3k=12$, so $(i,j,k)$ is $(1,6,2)$, $(3,3,3)$ or
$(5,0,4)$ and
\[
 \abs{B_0}=2^9\left(\frac{9!}{1!\,6!\,2!}+\frac{9!}{3!\,3!\,3!}+\frac{9!}{5!\,0!\,4!}\right)
 =1\,053\,696>2^{20}=\abs{\gppA\cap[0,M)} .
\]
No sum of two digits and no doubled digit exceeds $10$, so a midpoint relation in base
$11$ holds coordinatewise, and three points of a Euclidean sphere cannot have one as the
midpoint of the other two unless they coincide. Hence $B_0$ is progression-free, and
$\max B_0\le(11^9-1)/2<M/2$.

Set $D=B_0\cup(\gppA\cap[M,\infty))$. This remains progression-free. Two low terms cannot
reach a high term because $2\max B_0<M$. For a possible progression $a<b<c$ with $a$ low
and $b,c$ high, let $S=3^j$ be the largest ternary place of $c$, so that
$S\le c\le(3S-1)/2$. If $b<S$, then $2b\le S-1<c$. If $b\ge S$, then
$2b-c\ge(S+1)/2\ge(M+1)/2>a$. Both alternatives exclude $a+c=2b$. The difference
\[
 J(u)=\sum_{b\in B_0}u^b-\sum_{\substack{a\in\gppA\\a<M}}u^a
\]
is a polynomial with $J(1)=1\,053\,696-2^{20}>0$, hence positive at some real $z<1$, and
\cref{gpp:lem:replacement} applies with $L=M$. \gppsloppy The cardinality and the range are verified
by \artifact{grh-methods-polynomial-density/check_sphere.py}.
\end{example}

The failure of the conjecture does not depend on particular digit counts. It reflects a
growth gap between two digit constructions.

\begin{proposition}[Any sparse deficit rule can be improved; \stU]\label[proposition]{gpp:prop:growth}
Let $A_0\subseteq\Z_{\ge0}$ contain no nonconstant three-term arithmetic progression and
satisfy $\abs{A_0\cap[0,x]}\ll(1+x)^\alpha$ for some $0\le\alpha<1$. The deficit rule
$d\in A_0$ has limiting ball density $(1-q^{-1})\sum_{d\in A_0}q^{-d}$, and for every $q$
there is a polynomial-ratio progression-free set of strictly greater upper density. For
$A_0=\gppA$ one can take $\alpha=\log_32$.
\end{proposition}

\begin{proof}
For $X\ge1$, Abel summation gives $\sum_{d\in A_0}e^{-d/X}\ll_\alpha X^\alpha$, with a
constant depending also on the counting bound. Choose an integer $h\ge2$ so large that
$h>(2h-1)^\alpha$. In dimension $n$ read the digits $0,\ldots,h-1$ in base $2h-1$ and set
$X=(2h-1)^n$. One squared-norm level contains at least $h^n/(1+n(h-1)^2)$ vectors. Its
integer encodings form a progression-free set $B\subset[0,X/2]$ by the carry-free
midpoint argument of \cref{gpp:ex:sphere}. Consequently
\[
 \sum_{d\in B}e^{-d/X}\ge\frac{e^{-1/2}h^n}{1+n(h-1)^2}>\sum_{d\in A_0}e^{-d/X}
\]
for all sufficiently large $n$. Apply \cref{gpp:lem:replacement} with $D=B$,
$L=\max B$ and $z=e^{-1/X}$. For the ternary set there are $2^k$ permitted integers
below $3^k$, and already $h=4$ satisfies $4>7^{\log_32}$.
\end{proof}

\subsection{An explicit margin}\label{gpp:sec:margin}
{\raggedright\sources{grh_beyond_dirichlet_2026-11-25; grh-methods-polynomial-density}}

The proof of \cref{gpp:thm:counterexample} does not exhibit the first positive
coefficient $\Delta_R$. A variance bound for the degree of a random smooth polynomial
locates one in an explicit range and gives an explicit lower bound for $\beta_q-m_q$.

\begin{lemma}[Positive coefficient from a concentrated generating function; \stU]\label[lemma]{gpp:lem:concentration}
Let $H(u)=\sum_{n\ge0}H_nu^n$ have nonnegative coefficients and $0<H(z)<\infty$ for some
$0<z<1$. Define a random variable $X$ by $\gppPr(X=n)=H_nz^n/H(z)$. Suppose
$J(u)=P(u)-Q(u)$, where $P,Q$ have nonnegative coefficients, all exponents in $P$ are at
most $L$, and
\[
 J(z)>\gamma>0,\qquad P(z)\le C .
\]
If $\gppE X\le U$, $\gppVar X\le V$, and $t>0$ satisfies $CV/t^2<\gamma$, then for every
integer $R\ge L+U+t$,
\[
 \sum_{n=0}^R[u^n]\bigl(H(u)J(u)\bigr)\frac{z^n}{H(z)}>0 .
\]
In particular some coefficient of $HJ$ through $R$ is positive; if the coefficients are
integers, some coefficient is at least one.
\end{lemma}

\begin{proof}
The positive contribution omitted beyond $R$ is
\[
 \sum_{d\le L}[u^d]P(u)\,z^d\,\gppPr(X>R-d)\le P(z)\gppPr(X>U+t)\le\frac{CV}{t^2}
\]
by Chebyshev's inequality, since $\gppE X\le U$. Omitting the negative tail can only
increase the truncated signed sum. Its normalized value therefore exceeds
$\gamma-CV/t^2>0$.
\end{proof}

Take $z=127/128$, $L=480$ and the $74$-point set $B$ of \cref{gpp:thm:counterexample}.
Write $I_q(d)$ for the number of monic irreducibles of degree $d$ over $\Fq$, computed
recursively from $\sum_{e\mid d}eI_q(e)=q^d$. Under the law of
\cref{gpp:lem:concentration} for the smooth generating function $H$ of
\cref{gpp:lem:replacement}, the exponents of the different irreducibles are independent
geometric random variables with ratios $z^d$. Consequently
\[
 \gppE X=\sum_{d=1}^{480}I_q(d)\frac{dz^d}{1-z^d},\qquad
 \gppVar X=\sum_{d=1}^{480}I_q(d)\frac{d^2z^d}{(1-z^d)^2}.
\]
Define the integers
\begin{align}
 U_q^{\circ}&=\sum_{d=1}^{480}\ceil*{\frac{I_q(d)\,d\,127^d}{128^d-127^d}},\qquad
 V_q^{\circ}=\sum_{d=1}^{480}\ceil*{\frac{I_q(d)\,d^2\,127^d128^d}{(128^d-127^d)^2}},
 \notag\\
 R_q^*&=480+U_q^{\circ}+\ceil*{\sqrt{7400\,V_q^{\circ}}}. \label{gpp:eq:Rstar}
\end{align}
Termwise ceilings give $\gppE X\le U_q^{\circ}$ and $\gppVar X\le V_q^{\circ}$.

\begin{theorem}[Explicit search cutoff and margin; \stU]\label{gpp:thm:cutoff}
For the $74$-point witness some integer $R\le R_q^*$ satisfies $[u^R]H(u)J(u)\ge1$. In
particular, with $c_q$ as in \eqref{gpp:eq:delta} for $L=480$,
\begin{equation}\label{gpp:eq:explicitgap}
 \beta_q\ge m_q+(q-1)c_q\,q^{-R_q^*-1}>m_q .
\end{equation}
\end{theorem}

\begin{proof}
Use $P(u)=\sum_{b\in B}u^b$ and $Q(u)=\sum_{a\in\gppA}u^a$. By
\eqref{gpp:eq:Jpositive} one may take $\gamma=3/200$, and $P(z)\le74$. Put
$t=\ceil{\sqrt{7400V_q^{\circ}}}$. Then $74V_q^{\circ}/t^2\le1/100$, and the normalized
truncated sum is greater than $3/200-1/100=1/200>0$. \Cref{gpp:lem:concentration} provides
a positive integer coefficient with index at most $R_q^*$, and \eqref{gpp:eq:delta}
proves \eqref{gpp:eq:explicitgap}.
\end{proof}

\begin{corollary}[The binary margin; \stU]\label[corollary]{gpp:cor:binary}
One has $\beta_2>m_2+2^{-R_2^*-13}>m_2+2^{-2^{476}}$.
\end{corollary}

\begin{proof}
Exact integer arithmetic gives a value of $R_2^*$ of bit length $476$ with
$R_2^*+13<2^{476}$; the integer is recorded in
\artifact{grh_beyond_dirichlet_2026-11-25/data/density_certificate.json}. Using
$I_2(d)\le2^d/d$ and $\log(1+t)\le t$,
\[
 c_2=\frac12\prod_{d=1}^{480}(1+2^{-d})^{-I_2(d)}
 \ge\frac12\exp\Bigl(-\sum_{d=1}^{480}\frac1d\Bigr)\ge\frac1{960e}>2^{-12},
\]
since $\sum_{d\le480}1/d\le1+\log480$ and $e<3$. The additional density in
\eqref{gpp:eq:explicitgap} is $c_2\,2^{-R_2^*-1}>2^{-R_2^*-13}$.
\end{proof}

The certified cutoffs satisfy $R_q^*<2^{b}$ with
\[
\begin{array}{c|ccccc}
 q&2&3&4&5&7\\\hline
 b&476&756&956&1110&1343
\end{array}
\]
all computed with integers. The margin remains extraordinarily small: the proof shows
that a positive coefficient exists within an explicit range and does not enumerate an
extremal set.

\subsection{Finite optimization problems converging to the supremum}\label{gpp:sec:supremum}
{\raggedright\sources{grh-methods-polynomial-density}}

For the integers, McNew computes extremal upper densities by optimizing over a finite set
of primes and then enlarging it \cite[Theorems~4.2 and~5.1]{McNew2015Sets}. The following
adaptation to all irreducibles of degree at most $T$ bounds the approximation error in
terms of the degree cutoff.

\begin{theorem}[Computability and attainment of $\beta_q$; \stU]\label{gpp:thm:supremum}
Fix $q$. For an integer $T\ge1$ let $S$ be the set of all monic irreducibles of degree at
most $T$, and let $J_S(t)$ be the largest size of a polynomial-ratio progression-free
subset of the monic $S$-smooth polynomials of degree at most $t$. Put
\[
 U_T=(1-q^{-1})\prod_{P\in S}(1-q^{-\deg P})\sum_{t\ge0}J_S(t)q^{-t},\qquad
 v_T=\frac{1-q^{-1}}{\prod_{P\in S}(1-q^{-2\deg P})}.
\]
Then
\[
 v_TU_T\le\beta_q\le U_T,\qquad
 0\le U_T-\beta_q\le1-v_T\le\frac{q^{-T}}{q-1}.
\]
Thus $\beta_q=\lim_{T\to\infty}U_T$ is a computable real number. Moreover the supremum
$\beta_q$ is attained as the upper density of one progression-free set. The same
statements hold under the rational-ratio convention, if that convention is used both for
the supremum and for every $J_S(t)$.
\end{theorem}

\begin{proof}
For the upper bound, fix the scalar and the monic $S$-rough part $g$ of a polynomial. At
cutoff $N$, at most $J_S(N-\deg g)$ smooth factors can occur in that class. For the lower
bound, restrict the allowed $g$ to be squarefree and use a maximizing subset independently
in each class, with the same choice for all scalars at fixed $g$. In any allowed
progression the valuations outside $S$ are three equally spaced members of $\set{0,1}$,
hence constant. The three monic parts therefore lie in one class, where their quotients
would give a prohibited progression between smooth factors. This also handles signed
valuation steps in the rational-ratio version, including nonconstant ratios of degree
zero.

The proportion of monic polynomials of degree $n$ coprime to all primes in $S$ tends to
$c=\prod_{P\in S}(1-q^{-\deg P})$, by finite inclusion--exclusion. The proportion of those
that are also squarefree tends to $cv_T$, by the generating function in the proof of
\cref{gpp:lem:replacement}. In the normalized count the classes with $\deg g=N-t$
therefore contribute at most $(1-q^{-1})c\,J_S(t)q^{-t}$ in the limit, and the lower
construction attains the same expression times $v_T$. Dominated convergence applies since
$J_S(t)\le(t+1)^{\abs{S}}$, by counting the possible exponent vectors. Also $U_T\le1$, by
applying the identical count with all smooth factors allowed. Finally
\[
 1-v_T=1-\prod_{P\notin S}(1-q^{-2\deg P})\le\sum_{d>T}q^dq^{-2d}=\frac{q^{-T}}{q-1}.
\]
To pass from the finite constructions to one infinite set, use the degree blocks
\eqref{gpp:eq:blocks}. A rational-ratio progression with nonconstant degrees can be
ordered by increasing degree, by reversing its ratio if necessary, and triples of equal
degree are confined to one block.

For each fixed $T$, finite enumeration of admissible sets computes $J_S(t)$ exactly, and
truncating the series at $t=L$ costs at most the explicit tail
$\sum_{t>L}(t+1)^{\abs{S}}q^{-t}$. This determines $\beta_q$ to any prescribed error.

For attainment, at stage $j$ take $T=j$ and choose $N_j$ so large that the squarefree
construction has density at least $v_jU_j-1/j$ in its ball and that
\eqref{gpp:eq:blocks} holds. Keep its top $j+1$ degree layers. The union is
progression-free and has upper density at least
$\lim_{j\to\infty}(v_jU_j-1/j-q^{-j-1})=\beta_q$; equality follows from the definition of
$\beta_q$.
\end{proof}

\subsection{Upper bounds for polynomial ratios}\label{gpp:sec:upper}
{\raggedright\sources{grh-progression-densities-and-cubic-prime-race; mcnew-density-and-dirichlet-grh; grh-methods-polynomial-density}}

Let $r_3(n)$ be the largest size of a subset of $\set{0,\ldots,n-1}$ free of nonconstant
three-term arithmetic progressions, with $r_3(0)=0$. Write $z=1/q$. The one-prime upper
bound of \cite{Asada2017Subsets} is
\[
 U_q=(q-1)\sum_{n\ge1}\bigl(r_3(n)-r_3(n-1)\bigr)q^{-n}.
\]
A second linear prime forces strict losses against this bound, which answers
Question~6.2 of \cite{Asada2017Subsets}.

\begin{theorem}[Two-prime upper bound; \stU]\label{gpp:thm:twoprime}
For every prime power $q$,
\[
 \beta_q\le U_q-(1-z)^2(z^8+z^{19}+z^{29}).
\]
In particular $\beta_2\le907739729/1073741824<0.845398501494$.
\end{theorem}

\begin{proof}
Choose distinct monic linear primes $P,Q$ and a polynomial-ratio progression-free set.
Up to degree $N$, the nonzero polynomials of the form $bP^iQ^j$, for fixed $b$ coprime
to $PQ$ with its leading coefficient retained, have indices in
\[
 T_L=\set{(i,j)\in\Z_{\ge0}^2:i+j<L},\qquad L=N-\deg b+1 .
\]
In each row and each column the selected positions are progression-free in the varying
index, since a nonconstant progression on an axis gives a polynomial ratio. Summing the
row bounds permits $R_L=\sum_{k=1}^Lr_3(k)$ positions, and the same bound is obtained by
columns. Thus the total deficiency $R_L$ minus the number of selected positions can be
computed either by summing row deficiencies or by summing column deficiencies.

Suppose at least $h$ deficiencies are forced in $T_K$. For $L\ge K$, the last $K$ rows of
$T_L$ form a translated $T_K$, and so do the last $K$ columns; their lengths are unchanged
by translation. Consequently the deficiency of $T_L$ is at least $h$, and if it equals $h$
then every row and every column indexed below $L-K$ attains its maximum. For $K=0$ take
$h=0$. The following facts about intervals, with indices starting at zero, give three
contradictions to equality.
\[
\begin{array}{c|c|c|c|c}
 L&K&i&\text{row }i:\ \text{all maxima include }0&
         \text{column }0:\ \text{all maxima omit }i\\ \hline
 9&0&4 &r_3(5)=4&r_3(9)=5\\
 20&9&7&r_3(13)=7&r_3(20)=9\\
 30&20&4&r_3(26)=11&r_3(30)=12
\end{array}
\]
In each line $i<L-K$, so the two extremizers would disagree at the position $(i,0)$.
Starting with $T_0$ empty and applying the argument successively shows that the
deficiency is at least $1,2,3$ at the displayed lengths. Embedding these three tips in
longer triangles gives a penalty of
$\mathbf 1_{L\ge9}+\mathbf 1_{L\ge20}+\mathbf 1_{L\ge30}$ in each fibre.

\gppsloppy The interval entries are established by a finite enumeration. Generate the increasing
progression-free subsets by appending a permitted index $t$ and removing from the future
choices each $2t-a$ with $a$ already selected; this visits every progression-free subset
exactly once. Group by maximum and size, then take unions and intersections of the masks
of maximum cardinality. The program
\artifact{grh-progression-densities-and-cubic-prime-race/code/triangle_check.py} verifies
the table with integer arithmetic. For the first line one checks directly that the only
set of size four on five positions is $\set{0,1,3,4}$ and that the sets of size five on
nine positions are $\set{0,1,3,7,8}$, $\set{0,1,5,6,8}$, $\set{0,1,5,7,8}$ and
$\set{0,2,3,7,8}$.

To sum the row bounds, merge each row into a one-prime fibre with $P$ as the varying
prime and $P$-free base $bQ^j$. After division by $q^{N+1}$, the number of such bases of
degree $N-k$ tends to $(1-z)^2z^k$ at fixed $k$, by counting the polynomials of degree
$N-k$ and excluding the multiples of $P$. Their contribution is bounded in the limit by
\[
 (1-z)^2\sum_{k\ge0}r_3(k+1)z^k=U_q,
\]
where the equality follows by summing differences and domination follows from
$r_3(k+1)\le k+1$. Each penalty for a threshold $L\ge s$ is counted once for each $b$
coprime to $PQ$ of degree at most $N+1-s$; by inclusion--exclusion the normalized count
tends to $(1-z)^2z^{s-1}$. This proves the bound. For the numerical value the checker
computes $r_3(n)$ through $n=32$ and bounds the remaining increments by one, giving a
tail of at most $2^{-32}$ at $q=2$.
\end{proof}

The local calculation uses only ratios that are powers of one of two distinct linear
primes. Together with \cref{gpp:cor:binary},
\[
 m_2+2^{-2^{476}}<\beta_2\le\frac{907739729}{1073741824}<0.845398501494,
 \qquad m_2=0.8453979555\ldots,
\]
so $0<\beta_2-m_2<5.46\cdot10^{-7}$.

The next estimate locates the size of the discrepancy as the field grows.

\begin{theorem}[The supremum at large field size; \stU]\label{gpp:thm:largeq}
For every prime power $q\ge16$,
\[
 0<\beta_q-m_q\le(2q+1)q^{-q+14}.
\]
In particular $\beta_q$ and $m_q$ agree to every fixed order in $1/q$ as $q\to\infty$
through prime powers, and the upper limit of $\log(\beta_q-m_q)/(q\log q)$ is at most
$-1$.
\end{theorem}

\begin{proof}
The lower bound is \cref{gpp:thm:counterexample}. Put $K=q$, fix a degree cap $N\ge K$
and a progression-free set $E$. Call the set of polynomials of degree $N-k$ layer $k$,
and call a layer preferred if $k\in\gppA$. For each $k\in[0,K]\setminus\gppA$, form $a$
and $b$ by replacing every ternary digit $2$ of $k$ by $0$ and by $1$, respectively,
keeping the other digits. Then $a,b\in\gppA$, $a<b<k$ and $k+a=2b$. Set $d=k-b$. For
every $f\in E$ on layer $k$ and every monic irreducible $P$ of degree $d$, at least one
of $fP$ on layer $b$ and $fP^2$ on layer $a$ is missing from $E$. Assign weight
$1/I_q(d)$ to one such missing polynomial. The weights sent by each $f$ sum to one. All
three polynomials have the same leading coefficient, so the count also applies to sets
which distinguish associates.

For a nonzero polynomial $g$ let $\omega_d(g)$ be the number of its distinct monic
irreducible factors of degree $d$. For fixed $k$ and target $g$, at most $\omega_d(g)$
choices of $P$ can send weight to $g$; the layer of $g$ determines whether $f$ is
recovered as $g/P$ or as $g/P^2$. Moreover a target layer $h$ and a ratio degree $d$
determine at most one source layer $k$. Indeed the ternary digits of $d$ are ones
precisely where $k$ has twos. On this nonempty set of positions $b$ has only ones and $a$
has only zeros. Thus $h$ cannot occur in both roles for the same $d$, and the role fixes
$k$ as $h+d$ or $h+2d$.

Choose
\[
 \theta_1=1-\frac8q,\qquad\theta_d=\frac{8}{qd(d-1)}\quad(d\ge2),\qquad
 \sum_{d\ge1}\theta_d=1 .
\]
Call a pair $(k,g)$ exceptional if $\omega_d(g)>\theta_dI_q(d)$. Across all its
nonexceptional pairs, $g$ receives at most one unit of weight, by the uniqueness of the
source layer for each $d$. Hence the missing elements on preferred layers account for all
nonexceptional weight. The weight at an exceptional pair is at most one.

For $g$ uniform among the polynomials of a prescribed degree $n$, also with prescribed
leading coefficient, the probability that each of $t$ prescribed distinct primes of
degree $d$ divides $g$ is at most $q^{-dt}$. Taking $t=\floor{\theta_dI_q(d)}+1$ and
using the union bound gives
\[
 \gppPr\bigl(\omega_d(g)\ge t\bigr)\le\binom{I_q(d)}{t}q^{-dt}.
\]
For $d=1$ we have $I_q(1)=q$ and $t=q-7$, and the probability is at most
$\binom q7q^{-(q-7)}\le q^{-q+14}$. For $d\ge2$,
$dI_q(d)\ge q^d-\sum_{j=1}^{\floor{d/2}}q^j\ge q^d/2$, since every element of
$\F_{q^d}$ of degree less than $d$ over $\Fq$ lies in a proper subfield of degree at most
$d/2$. Hence
\[
 t>\frac{4q^{d-1}}{d^2(d-1)}\ge q,\qquad
 \frac{e}{\theta_dq^d}=\frac{e\,d(d-1)}{8q^{d-1}}\le\frac{e}{4q}<\frac1q .
\]
Both bounds hold at $d=2$; when $d$ increases by one, the denominators $d^2(d-1)$ and
$d(d-1)$ grow by factors at most $9/2$ and $3$, respectively, while $q\ge16$. Using
$\binom{I}{t}\le(eI/t)^t$ now bounds the probability by $q^{-q}$ for $d\ge2$.

Write $B_N^\times$ for the set of nonzero polynomials of degree at most $N$. Across all
$k\le K$ the exceptional weight is at most $2q\abs{B_N^\times}q^{-q+14}$. The proportion
of $B_N^\times$ on layers deeper than $K$ is $(q^{N-K}-1)/(q^{N+1}-1)\le q^{-K-1}$. The
count of $E$ is at most the full count on preferred layers, plus the exceptional weight
and this allowance for deeper layers: every remaining member on a nonpreferred layer has
one unit of weight, which has been charged against a missing element of a preferred
layer. The proportion of $B_N^\times$ on preferred layers tends to
$(1-q^{-1})\sum_{j\in\gppA}q^{-j}=m_q$. Taking the upper limit in $N$, then the supremum
over $E$, and using $K=q$ proves the bound.
\end{proof}

\subsection{Natural density, and rational-function ratios}\label{gpp:sec:natural}
{\raggedright\sources{grh-progression-densities-and-cubic-prime-race}}

The benchmark for natural density is the greedy set $\gppG$ of
\cref{gpp:sec:conventions}. A finite change in the allowed joint valuations at a few
primes gives a larger natural density, which answers Question~6.1 of
\cite{Asada2017Subsets}; the precedent is McNew's exchange at $2,3,5$ over the integers
\cite[Theorem~3.3]{McNew2015Sets}.

Let $S$ be a finite nonempty set of distinct monic irreducibles and put $k=\abs{S}$.
Retain the constraint $v_P(f)\in\gppA$ at every $P\notin S$. For the coordinates indexed
by $S$, replace $\gppA^k$ by
\[
 D_S=\set{u\in\gppA^k:\min_iu_i=0}\ \cup\
 \bigl(\set{1,2}^k\setminus\set{(1,\ldots,1)}\bigr),
\]
and call the resulting set of polynomials $E_S$. Deleting the entire positive part of
$\gppA^k$ makes it possible to account for every higher-power conflict in one step.

\begin{lemma}[\stU]\label[lemma]{gpp:lem:ES}
The set $E_S$ contains no triple with a nonconstant polynomial ratio.
\end{lemma}

\begin{proof}
The common valuation step of a polynomial ratio is nonnegative at each prime. All its
coordinates outside $S$ vanish since $\gppA$ is progression-free. Inside $S$ write the
three vectors as $u,u+h,u+2h$ with $h\ge0$, $h\ne0$. If the final point has a zero
coordinate, all three points have a zero and belong to $\gppA^k$, which is impossible.
Otherwise the final point is in the new cube. At every moving coordinate its final value
is $2$, its initial value $0$ and its middle value $1$. The initial point therefore
belongs to $\gppA^k$. Any fixed coordinate must be $1$, since its value lies in both
$\gppA$ and $\set{1,2}$. The middle point is then the deleted all-ones vector.
\end{proof}

\begin{theorem}[An explicit exchange beating the greedy set in natural density; \stU]\label{gpp:thm:natural}
For every prime power $q$ an explicit finite set $S$ gives
$\gppdens(E_S)>\gppdens(\gppG)$. With $T(w)=\sum_{a\in\gppA,\ a\ge1}w^{a-1}$ and the
notation of \cref{gpp:lem:valuations}, the exact gain is
\begin{equation}\label{gpp:eq:gain}
 \gppdens(E_S)-\gppdens(\gppG)=C_S\prod_{P\in S}(1-w_P)w_P
 \left[\prod_{P\in S}(1+w_P)-1-\prod_{P\in S}T(w_P)\right].
\end{equation}
One may take all $q$ monic linear polynomials when $q\ge4$, the set
$\set{x,x+1,x^2+x+1,x^3+x+1}$ when $q=2$, and the three monic linear polynomials together
with $x^2+1$ when $q=3$.
\end{theorem}

\begin{proof}
Both densities exist by \cref{gpp:lem:valuations}. The two terms subtracted and added
inside the brackets of \eqref{gpp:eq:gain} are the weights of the positive cube with its
all-ones point omitted and of the old positive orthant, after division by the weight of
the all-ones point.

For $q\ge4$ and all $q$ monic linears: the first two positive members of $\gppA$ are $1$
and $3$, so $T(w)\le1+w^2/(1-w)$ and
\[
 T(1/q)^q\le e^{1/(q-1)}\le e^{1/3}<1.4,\qquad
 (1+1/q)^q-1\ge(5/4)^4-1>1.44 ;
\]
the lower bound follows by expanding the binomial, since each term of fixed degree in
$(1+1/q)^q$ increases with $q$. For $q=2,3$ use the sharper bound
$T(w)\le1+w^2+w^3+w^8/(1-w)$, which follows from $\gppA\cap[1,8]=\set{1,3,4}$, the bound
\eqref{gpp:eq:CS}, and the degree multisets in the table. Substitution in
\eqref{gpp:eq:gain} gives the strict lower bounds shown, rounded down; the line for $q=4$
illustrates the size of the improvement.
\begin{center}
\begin{tabular}{@{}clll@{}}\toprule
$q$ & Degrees in $S$ & Bracket in \eqref{gpp:eq:gain} exceeds & Density gain exceeds\\\midrule
2 & $1,1,2,3$ & $0.06622$ & $0.00004244$\\
3 & $1,1,1,2$ & $0.09852$ & $0.00007119$\\
4 & $1,1,1,1$ & $0.09023$ & $0.00008364$\\\bottomrule
\end{tabular}
\end{center}
\gppsloppy The displayed quadratic and cubic polynomials have no root in their ground fields. The
table is reproduced with rational arithmetic by
\artifact{grh-progression-densities-and-cubic-prime-race/code/check_density.py}.
\end{proof}

\begin{remark}\label[remark]{gpp:rem:mixed}
The monotonicity of the valuation step is essential in \cref{gpp:lem:ES}. For two
selected coordinates the three points $(0,3),(1,2),(2,1)$ all obey the new local rule and
have the mixed common step $(1,-1)$; in larger dimension append coordinates constantly
equal to $1$. Thus $E_S$ is not free of rational-ratio triples, and that convention needs
separate constructions: \cref{gpp:thm:rational} below, and the one-digit exchange of
\cref{gpi:prop:poly-greedy}.
\end{remark}

The bare existence of a polynomial-ratio progression-free set of natural density greater
than $\gppdens(\gppG)$ also follows from \cref{gpi:prop:poly-greedy}, which produces such a
set under the stronger rational-ratio convention; every rational-ratio progression-free
set is polynomial-ratio progression-free. That construction needs a linear marker and
satellite primes of total weight exceeding $q^2+q-1$ and comes with no numerical gain.
\Cref{gpp:thm:natural} is retained for what it adds under the polynomial-ratio convention
of Question~6.1: the exchange uses only $\max(4,q)$ primes, its gain is given exactly by
\eqref{gpp:eq:gain}, and the table gives certified numerical gains. Neither result
implies the other.

For the rational-ratio convention fix a monic linear $P$. Every nonzero polynomial has a
unique expression $cgP^e$ with $c\in\Fq^\times$, $g$ monic, $P\nmid g$ and $e\ge0$. Call
$g$ \emph{admissible} when all its valuations belong to $\gppA$, and let $g_n$ be the
number of admissible cofactors of degree $n$. \Cref{gpp:lem:valuations} gives
\[
 g_n/q^n\longrightarrow c_P=(1-q^{-1})\prod_{Q\ne P}p_Q>0 .
\]
For every admissible $g$ choose a set $I_g\subset\Z_{\ge0}$ free of nonconstant
three-term arithmetic progressions, and select exactly the polynomials $cgP^e$ with $g$
admissible and $e\in I_g$. Every such set is free of rational-ratio triples. Indeed a
rational-function ratio has signed valuation steps; the restrictions outside $P$ force
all of those steps to be zero, so the three terms have the same monic cofactor $g$, and a
nonconstant ratio would have to change the $P$-exponent by a nonzero step.

\begin{theorem}[Rational-function ratios with one free linear prime; \stU]\label{gpp:thm:rational}
Let $s(t)=r_3(t+1)$ be the largest size of a three-term-progression-free subset of
$\set{0,\ldots,t}$. In the family just described, the maximum attainable upper density is
\[
 c_P(1-q^{-1})\sum_{t\ge0}q^{-t}s(t).
\]
This maximum exceeds $\gppdens(\gppG)$ by at least
$c_P(1-q^{-1})q^{-8}\ge(1-q^{-1})^3q^{-8}$, although the greedy restriction is relaxed at
one linear prime only.
\end{theorem}

\begin{proof}
At cutoff $N$ any admissible $g$ of degree $N-t$ permits at most $s(t)$ exponents. The
normalized count is at most $(q-1)q^{-N-1}\sum_{t\le N}g_{N-t}s(t)$. For fixed $t$ the
summand with its factor tends to $c_P(1-q^{-1})q^{-t}s(t)$, and the terms are dominated
by $(1-q^{-1})q^{-t}(t+1)$, which is summable. This proves the upper bound.

To realize it in a single set, choose $N_j=2^{j+4}$ for $j\ge1$ and $N_0=-1$. For each
cofactor with $N_{j-1}<\deg g\le N_j$, let $I_g$ be a progression-free maximizer in
$[0,N_j-\deg g]$, with ties broken lexicographically. This defines each fibre once. At
the cutoff $N_j$ the current block attains the upper bound term by term for
$0\le t<N_j-N_{j-1}$, and the bound on the omitted normalized tail tends to zero.
Therefore the upper density equals the claimed maximum.

For the greedy set the same counting identity holds with $\abs{\gppA\cap[0,t]}$ in place
of $s(t)$, since every fibre is $\gppA$. All differences are nonnegative. At $t=8$ the
set $\gppA$ selects only $0,1,3,4$, whereas $\set{0,1,3,7,8}$ is progression-free of size
five. This single summand gives the stated gain, and $\prod_{Q\ne P}p_Q\ge1-q^{-1}$ by
\eqref{gpp:eq:CS}.
\end{proof}

For $q=2$ the guaranteed gain is $1/2048$. The limits in \cref{gpp:thm:rational} are
upper densities along cutoffs, whereas \cref{gpp:thm:natural} gives a natural density.
In particular the existence clause of Question~6.4 of \cite{Asada2017Subsets} has a
positive answer. That clause, taken alone, is also a consequence of
\cref{gpi:prop:poly-greedy}, which gives a rational-ratio progression-free set whose
natural density exceeds $\gppdens(\gppG)$, and of \cref{gpi:thm:poly-gap}, which bounds
the supremum $u_q$ of the upper densities of such sets from below by the supremum $b_q$ of
their natural densities plus an explicit positive quantity, greater than $0.003660$ for
$q=2$; since $b_q\ge\gppdens(\gppG)$ this is a larger excess than $1/2048$.
\Cref{gpp:thm:rational} is not implied by those results: it determines the exact optimum
of the family in which the greedy condition is retained outside one linear prime. By
\cref{gpp:thm:supremum} the unrestricted rational-ratio supremum $u_q$ is the limit of the
finite optimization problems $U_T$ formed with the rational-ratio convention; the bounds
$v_TU_T\le u_q\le U_T$ are also proved in \cref{gpi:thm:poly-gap}. For polynomial ratios,
the supremum of the natural densities is strictly smaller than $\beta_q$
(\cref{gpi:rem:poly-ratio}).

\subsection{Finite coefficient rings: local rings with square-zero maximal ideal and their products}\label{gpp:sec:rings}
{\raggedright\sources{grh_method_transfers; grh_method_transfers_graph_towers_and_finite_rings; grh_method_transfers_research}}

Question~6.5 of \cite{Asada2017Subsets} asks how the density problem changes when another
finite ring replaces the coefficient field. The nonunit-ratio convention of that paper is
retained: ratios are polynomials that are not units. A finite product of rings shows a new
feature, since multiplication by $x$ can be switched on independently in any subset of the
factors; nilpotents show another, since the ratio and its square then act differently.

All rings are commutative with identity, and local rings have a nonzero residue field. For
a finite ring $R$ put $B_n(R)=\set{f\in R[x]:\deg f\le n}$, including zero, with
$\deg0=-\infty$; thus $\abs{B_n(R)}=\abs{R}^{n+1}$. For an integer $L\ge2$ call a set
\emph{admissible} if it contains no three distinct polynomials of the form
\begin{equation}\label{gpp:eq:pattern}
 f,\quad r^af,\quad r^{a+b}f,\qquad
 r\in R[x]\text{ a nonunit},\quad a,b\ge1,\quad a+b\le L .
\end{equation}
Thus $L=2$ is the progression problem of \cite{Asada2017Subsets}. Write $\Delta_L(R)$ for
the supremum over admissible $E$ of $\limsup_{n\to\infty}\abs{E\cap B_n(R)}/\abs{R}^{n+1}$,
and $\Delta(R)=\Delta_2(R)$. The proofs work also if triples with repeated terms are
forbidden: in the lower-bound constructions polynomials with a zero residue component are
discarded, and all triples used in upper bounds are distinct.

\begin{theorem}[Criterion for growing products; \stU]\label{gpp:thm:local}
Let $(A_i,\gppm_i)$ be finite local rings with $\gppm_i^2=0$. Write
$q_i=\abs{A_i/\gppm_i}$, $Q_i=\abs{A_i}$ and $R_s=\prod_{i=1}^sA_i$. Define
\[
 w_i=\begin{cases}Q_i^{-2},&q_i\text{ even},\\(q_iQ_i)^{-1},&q_i\text{ odd}.\end{cases}
\]
For every fixed $L\ge2$ the limit $\lim_{s\to\infty}\Delta_L(R_s)$ exists, and it is
positive if and only if $\sum_iw_i<\infty$. For one local factor, uniformly in $L$,
\[
 \gamma_iw_i\le1-\Delta_L(A_i)\le w_i,\qquad
 \gamma_i=\begin{cases}1-Q_i^{-1},&q_i\text{ even},\\(1-q_i^{-1})^2,&q_i\text{ odd}.\end{cases}
\]
For every finite $s$,
\begin{equation}\label{gpp:eq:productbound}
 \prod_{i=1}^s(1-w_i)\le\Delta_L(R_s)\le\gppE\,\alpha_K,\qquad
 K=\sum_{i=1}^sX_i,
\end{equation}
where the $X_i$ are independent Bernoulli variables with $\gppPr(X_i=1)=3\gamma_iw_i$ and
$\alpha_k$ is the largest density of a subset of the ternary word cube $\set{0,1,2}^k$
containing no combinatorial line, with $\alpha_0=1$.
\end{theorem}

The quantities $Q_i$ and $q_i$ measure different features. At fixed odd $q_i=q$ the
density loss has order $Q_i^{-1}$, whereas at fixed even $q$ it has order $Q_i^{-2}$. For
example take $A_i=\Fq\oplus V_i$ with $V_i^2=0$ and
$1+\dim V_i=\max(2,\ceil{\log_q(i+1)})$. Then $Q_i$ is comparable to $i+1$ for fixed $q$,
and the limiting extremal density is positive for $q=2$ and zero for $q=3$. The ring
$\Z/p^2\Z$ is covered with $q=p$ and $Q=p^2$, and a finite field is covered with $Q=q$
and $w=q^{-2}$, so that for products of fields the criterion is the convergence of
$\sum_iq_i^{-2}$; in particular $\Delta(\Fq^b)\to0$ as $b\to\infty$ for fixed $q$.

\begin{proof}
\emph{Degree restrictions at a single factor.} Let $(A,\gppm)$ be one of the factors,
$k=A/\gppm$, $q=\abs{k}$, $Q=\abs{A}$, and write $\bar f$ for coefficientwise reduction.
A polynomial is a unit if and only if its reduction is a nonzero constant: lift that
constant to a unit and invert the remaining nilpotent correction. A polynomial with
reduction of degree $e\ge1$ can be written $r=t+u$ with $\deg t=e$, the leading
coefficient of $t$ a unit, and $u\in\gppm[x]$; choose lifts only of the coefficients of
$\bar r$ up to degree $e$. Then
\[
 r^j=t^j+jt^{j-1}u\qquad(j\ge1).
\]
If the residue characteristic $p$ divides $j$, then $ju=0$, since $p\in\gppm$.
Multiplication by $t$ increases the degree of every nonzero polynomial by exactly $e$.
These assertions allow the characteristic of $A$ itself to be $p^2$.

Suppose $f$, $y=r^af$, $z=r^{a+b}f$ lie in $B_n(A)$ with all three reductions nonzero. If
$r$ is a nonunit, its reduction has positive degree $e$. In residue characteristic two at
least one of $a,b,a+b$ is even, and the corresponding relation $y=t^af$, $z=t^by$ or
$z=t^{a+b}f$ forces $f$ or $y$ to have degree at most $n-2$. Remove $B_{n-2}(A)$, of
relative size $Q^{-2}$. For odd $p$ remove instead
\[
 D_n=\set{f:\deg f\le n-1,\ \deg\bar f\le n-2},
\]
of relative size $(qQ)^{-1}$ by the independent uniform coefficients of $x^n$ and
$x^{n-1}$. If $p\mid b$, then $z=t^by$ forces $\deg y\le n-2$. Otherwise $b$ is a unit in
$A$, and elimination of $u$ gives
\begin{equation}\label{gpp:eq:eliminate}
 (a+b)t^by-az=bt^{a+b}f .
\end{equation}
Consequently $\deg f\le n-ae\le n-1$, and $\bar z=\bar r^{a+b}\bar f$ forces
$\deg\bar f\le n-(a+b)e\le n-2$, so the first term lies in $D_n$. After these deletions,
and the removal of the fibre of zero reduction at a cost of at most $q^{-(n+1)}$, no
pattern \eqref{gpp:eq:pattern} remains in $B_n(A)$.

\emph{Disjoint triples.} A \emph{packing} is a collection of vertex-disjoint ordered
triples \eqref{gpp:eq:pattern} with $a=b=1$; it is enough to specify the first terms and
an injective rule for each of the three positions. For even $q$ take every $f$ of exact
degree $n-2$, with triple $(f,xf,x^2f)$. The three positions have different degrees, and
the proportion of first terms is $(1-Q^{-1})Q^{-2}$. For odd $q$ and $n\ge3$ take all
$f=\sum f_jx^j$ with
\[
 f_n=0,\qquad f_{n-1}\in\gppm,\qquad f_{n-2},\,f_0\in A^\times .
\]
They form the proportion $(1-q^{-1})^2/(qQ)$ of the ball. Set
\[
 c=-\frac{f_{n-1}}{2f_{n-2}}\in\gppm,\qquad r=x+c(x^2+1),\qquad y=rf,\quad z=r^2f .
\]
The term $x^2c$ cancels the unwanted term of degree $n+1$ of $z$: indeed $cf_{n-1}=0$
and $f_{n-1}+2cf_{n-2}=0$, so $y,z\in B_n(A)$. Their residue degrees are $n-1$ and $n$,
distinct from the residue degree $n-2$ of the start. To recover the start, the
polynomial $\bar f$ is obtained from $\bar y=x\bar f$ or $\bar z=x^2\bar f$, and the
coefficients
\[
 [x^n]y=-cf_{n-2},\qquad[x]z=2cf_0
\]
determine $c$ from either position, because $\gppm$ is a vector space over $k$ and the
indicated scalars are known and nonzero in $k$. Finally $r(x-c(x^2+1))=x^2$ shows that
multiplication by $r$ and by $r^2$ is injective. The triples are pairwise disjoint. With
the deletions above this proves the bounds for one factor.

\emph{From finite balls to one infinite set.} For a fixed finite product let $M(n)$ be the
largest fraction of its ball occupied by an admissible set. Then
\begin{equation}\label{gpp:eq:glueproduct}
 \Delta_L(R_s)=\limsup_{n\to\infty}M(n).
\end{equation}
For the nontrivial inequality choose a subsequence $n_j$ approaching the upper limit so
rapidly that
\[
 n_j-j>(1+L)n_{j-1},\qquad n_j-(1+L)j>n_{j-1}.
\]
From an optimizer at $n_j$ keep only the elements whose residue degrees in all
coordinates lie in $[n_j-j,n_j]$. The fraction discarded is at most
$\sum_{i\le s}q_i^{-j-1}$, by uniform reduction. The union of the retained blocks is
admissible. In each coordinate the three residue degrees of a pattern are
$d,d+ae,d+(a+b)e$ with $e\ge0$, and all are finite. If only the last term were in the
latest block, its degree in each coordinate would be at most $(1+L)n_{j-1}$, because the
middle degree is at most $n_{j-1}$. If the last two were in that block, their difference
$be$ would be at most $j$, so the first degree would be at least $n_j-(1+L)j$. Both
contradict the separation. Slicing the product by its last coordinate shows
$\Delta_L(R_{s+1})\le\Delta_L(R_s)$, using ratio $1$ in the coordinate held fixed; hence
the limit in $s$ exists.

\emph{Products.} The product of the single-factor constructions, trimmed to nonzero
residues, gives the lower bound in \eqref{gpp:eq:productbound}: in a pattern the ratio
is a nonunit in at least one coordinate, where all three residue degrees are finite, so
the deletion argument in that coordinate applies. For the upper bound put
$u_i=\gamma_iw_i$ and partition $B_n(A_i)$ into the packed triples above and singleton
cells. A uniformly drawn point of the product ball falls in a triple with probability
$3u_i$ in coordinate $i$, independently across the coordinates. Conditionally on the
cells, if $K$ coordinates fall in triples the point is uniform on a ternary cube of
dimension $K$. Every combinatorial line gives a forbidden triple, since its moving
coordinates have their respective fixed ratios and the other coordinates can use ratio
$1$. Hence $M(n)\le\gppE\,\alpha_K$. The density Hales--Jewett theorem gives
$\alpha_k\to0$ \cite{Polymath2012New}. Since $\gamma_i\ge4/9$ and
$\gppVar K\le\gppE K$, divergence of $\sum_iw_i$ gives $K\to\infty$ in probability, and
the upper bound tends to zero. When the sum converges, the product in
\eqref{gpp:eq:productbound} is positive, since $w_i\le1/4$.
\end{proof}

\begin{proposition}[Products of finite fields; \stU]\label[proposition]{gpp:prop:fields}
For $R_b=\prod_{i=1}^b\F_{q_i}$ the upper bound in \eqref{gpp:eq:productbound} holds for
$\Delta(R_b)$ with the larger parameters
\[
 \gppPr(X_i=1)=\frac{3}{q_i^2+q_i+1}.
\]
For any sequence of finite products of fields, with no nesting assumption, set
$S=\sum_iq_i^{-2}$ for each product. Along such a sequence $\Delta\to0$ exactly when
$S\to\infty$, and $\Delta\to1$ exactly when $S\to0$; in the latter regime
\[
 1-\Delta=S+O(S^{3/2})
\]
uniformly.
\end{proposition}

\begin{proof}
Fix a field of order $q$. Each nonzero polynomial in its ball has a unique expression
$x^eg$ with $g(0)\ne0$. For each fixed $g$ partition the exponent interval
$0\le e\le n-\deg g$ into successive triples starting at $0,3,6,\ldots$, leaving up to
two singleton cells, and put zero in a singleton. For $3j+2\le n$ the number of triples
starting at exponent $3j$ equals $(q-1)q^{n-2-3j}$, the number of $g$ with nonzero
constant term and degree at most $n-2-3j$. The fraction of the ball lying in triple cells
is therefore
\[
 \theta_{q,n}=3(q-1)\sum_{0\le j\le(n-2)/3}q^{-3j-3}\longrightarrow\frac{3}{q^2+q+1}.
\]
The word-cube argument in the proof of \cref{gpp:thm:local} applies to these cells with
$r_i=x$ on the moving positions and $r_i=1$ elsewhere, and gives $\gppE\,\alpha_{K_n}$ as
an upper bound for the maximum in the finite ball; passage to the limit at fixed $b$
gives the first assertion. It uses only the $2^b-1$ nonconstant ratios whose coordinates
are in $\set{1,x}$.

Write $p_i=3/(q_i^2+q_i+1)$; these lie between $(12/7)q_i^{-2}$ and $3q_i^{-2}$. When
$S\to\infty$ the mean $\mu=\sum_ip_i$ of $K$ tends to infinity, $\gppVar K\le\mu$ gives
$\gppPr(K<\mu/2)\le4/\mu$, and the contribution of larger $K$ is bounded by
$\sup_{k\ge\mu/2}\alpha_k\to0$. Conversely the lower bound in
\eqref{gpp:eq:productbound} is at least $\exp(-4S/3)$, since $-\log(1-u)\le4u/3$ for
$0\le u\le1/4$. For the other endpoint, slicing a ternary cube of positive dimension into
parallel lines shows $\alpha_k\le2/3$ for $k\ge1$, whence
\[
 \frac{1-\prod_i(1-p_i)}{3}\le1-\Delta\le S .
\]
As $\sum_ip_i\ge12S/7$, the left side stays positive when $S$ stays bounded below. As
$S\to0$ each $q_i^{-1}\le\sqrt S$, so $\sum_ip_i=3S+O(S^{3/2})$, and the first two
inclusion--exclusion bounds give
$1-\prod_i(1-p_i)=\sum_ip_i+O((\sum_ip_i)^2)$. This proves the expansion.
\end{proof}

\begin{example}\label[example]{gpp:ex:fields}
With the values $\alpha_1=\alpha_2=\alpha_3=2/3$ from \cite{Polymath2010Density},
\cref{gpp:prop:fields,gpp:thm:local} give
\[
 \frac9{16}\le\Delta(\F_2\times\F_2)\le\frac{38}{49},\qquad
 \frac{16}{25}\le\Delta(\F_2\times\F_3\times\F_5)\le\frac{322}{403}.
\]
\gppsloppy The chain counts are verified by direct enumeration at small cutoffs, all $37$
lines in a three-coordinate cell are checked, and the displayed bounds are computed in
rational arithmetic, by the program \code{check_products.py} in
\artifact{grh_method_transfers_graph_towers_and_finite_rings/verification/}.
\end{example}

The following lower bound is valid for every finite product of local rings, with radicals
of arbitrary nilpotence order. The choices of $a,b$ in \eqref{gpp:eq:pattern} give
precisely the several weights that \cref{gpp:lem:digits} handles simultaneously.

\begin{proposition}[A Behrend-type lower bound; \stU]\label[proposition]{gpp:prop:behrendrings}
Let $R=\prod_{i=1}^sA_i$ be a finite product of finite local rings with residue fields of
orders $q_i$, and put $V=\sum_iq_i/(q_i-1)^2\le2s$. For a constant $C_L$ depending only on
$L$,
\[
 \Delta_L(R)\ge\exp\bigl(-C_L\sqrt{\log(2+V)}\bigr).
\]
\end{proposition}

\begin{proof}
At cutoff $n$ consider, on tuples with every residue component nonzero, the sum of the
residue degree deficits $D_n=\sum_i(n-\deg\bar f_i)$. For fixed $s$ its limiting law is
that of a sum $D$ of independent variables with masses $(1-q_i^{-1})q_i^{-j}$ for
$j\ge0$, of variance $V$; the stated geometric probability for a given $j\le n$ is the
exact probability of residue degree $n-j$, and zero components have vanishing
probability. By Chebyshev's inequality an interval $I$ of $m\le4\sqrt V+4$ consecutive
integers, enlarged to length two if needed, carries at least $3/4$ of the mass of $D$.
Let $S\subset\set{0,\ldots,m-1}$ be the set of \cref{gpp:cor:behrend}. As $h$ ranges over
the $2m-1$ integers of $I-\set{0,\ldots,m-1}$, each point of $I$ belongs to exactly
$\abs{S}$ translates $h+S$. Hence some translate receives probability at least
$3\abs{S}/(4(2m-1))$. At cutoff $n$ select the tuples with nonzero residue components and
$D_n\in h+S$. Along a pattern \eqref{gpp:eq:pattern} the sum of the residue degrees
increases by $a\sigma$ and then by $b\sigma$, where $\sigma=\sum_i\deg\bar r_i>0$, so the
three deficits would satisfy a forbidden weighted mean. Passing to the limit and applying
the block argument \eqref{gpp:eq:glueproduct}, which only uses reduction to the residue
fields, proves the bound.
\end{proof}

\begin{theorem}[The leading constant in odd residue characteristic; \stU]\label{gpp:thm:oddleading}
Let $A$ be a finite local ring with residue field $k=\Fq$ of odd characteristic, maximal
ideal $J$ with $J^2=0$, cardinality $Q$, and $M=\abs{J}=Q/q$. For every $L\ge2$,
\[
 \frac{1-q^{-q}}{qQ}\left(1-\frac{2+2q}{M}\right)
 \le1-\Delta_L(A)\le\frac{1-q^{-q}}{qQ}+\frac{q+1}{Q^2}.
\]
Consequently, along any sequence of such rings with fixed residue field of odd order $q$,
\[
 \lim_{Q\to\infty}qQ\bigl(1-\Delta_L(A)\bigr)=1-q^{-q}.
\]
\end{theorem}

The constant accounts for the proportion $q^{-q}$ of residue polynomials divisible by
every monic linear polynomial. With \cref{gpp:thm:local}, for ordinary three-term
progressions and fixed residue field,
\[
 1-\Delta(A)=\begin{cases}Q^{-2}+O(Q^{-3}),&q\text{ even},\\
 (1-q^{-q})/(qQ)+O_q(Q^{-2}),&q\text{ odd}.\end{cases}
\]

\begin{proof}
Multiplication makes $J$ a vector space over $k$. Choose a representative
$\tilde\gamma\in A$ of each $\gamma\in k$, and put $P_\gamma=x-\tilde\gamma$ and
$\Lambda=\prod_{\gamma\in k}(x-\gamma)\in k[x]$. Subscripts denote coefficients.

\emph{Construction.} At cutoff $n\ge q+3$ consider the rectangle $f_n=0$, $f_{n-1}\in J$,
of proportion $1/(qQ)$. Inside this rectangle delete only the points with
$\Lambda\nmid\bar f$; their proportion in the full ball is $(1-q^{-q})/(qQ)$ by
uniformity modulo the polynomial $\Lambda$ of degree $q$. Also delete all points of
degree at most $n-2$ and all points divisible in $A[x]$ by some $P_\gamma^2$. The
additional cost is at most $(q+1)Q^{-2}$, since $P_\gamma$ is monic. The points of zero
residue are removed at a cost tending to zero.

Consider a remaining pattern $f$, $y=r^if$, $z=r^{i+j}f$ with a nonunit ratio. Nonzero
reductions of all three terms ensure that $e_0=\deg\bar r$ is positive. Let $t$ be a
coefficient lift of $\bar r$ of degree $e_0$ with unit leading coefficient. If the
residue characteristic divides $j$, then $r^j=t^j$ and $y$ has degree at most
$n-je_0\le n-2$, which was excluded. Otherwise $j$ is a unit, and
\eqref{gpp:eq:eliminate} gives $\deg f\le n-ie_0$. Since the points of degree at most
$n-2$ were deleted, $i=e_0=1$ and $\deg f=n-1$; moreover
$\deg\bar f\le n-(i+j)e_0\le n-2$. Thus $f$ lies in the rectangle, and
$\Lambda\mid\bar f$ because $f$ survived.

Write $r=uP_\gamma+\epsilon$ with $u\in A^\times$ constant and $\epsilon\in J[x]$. At a
surviving first term $\Lambda\mid\bar f$, hence $f=P_\gamma g+\eta$ for some $g\in A[x]$
and $\eta\in J[x]$, and $\epsilon f=P_\gamma\epsilon g$. For $m=i+j\ge2$,
\[
 r^mf=u^mP_\gamma^mf+mu^{m-1}P_\gamma^{m-1}\epsilon f
\]
is divisible by $P_\gamma^2$, so the last term was deleted. Gluing the cutoffs by
\eqref{gpp:eq:glueproduct} gives the upper bound for $1-\Delta_L(A)$.

\emph{A matching of triples.} Start with all $f$ in the rectangle with
$\Lambda\nmid a$, where $a=\bar f$. Choose the first $\beta\in k$ in a fixed order with
$a(\beta)\ne0$. Put $P=P_\beta$, $s=n-\deg a\ge2$ and
\[
 c=\frac{-f_{n-1}}{2\gpplc(a)}\in J,\qquad r=P+c(P^s+1),\qquad y=rf,\quad z=r^2f,
\]
where the operations on $c$ use the $k$-action on $J$. These formulas give
$\deg y,\deg z\le n$ and $y_n=-c\,\gpplc(a)$: the coefficient of degree $n+1$ in
$z=P^2f+2cP(P^s+1)f$ is $f_{n-1}+2c\,\gpplc(a)=0$. Discard $c=0$. Each remaining first
term has degree $n-1$ and $y$ has degree $n$. Require $z_n\ne0$ as well. If $s=2$ the
reduction $\bar z=(x-\beta)^2a$ already has degree $n$. For $s>2$, conditionally on $a$
and $c$, the coefficient $f_{n-2}$ is uniform on $J$ and occurs in $z_n$ with coefficient
one, while its occurrences in terms multiplied by $c$ depend only on its reduction. Thus
these two exclusions cost at most the fraction $2/M$ of the starts.

At $\tilde\beta$ one has $y(\tilde\beta)=ca(\beta)\ne0$ and $z(\tilde\beta)=0$. If
$\gamma\ne\beta$ is a root of $a$, then $f(\tilde\gamma)\in J$ and
\[
 y(\tilde\gamma)=P(\tilde\gamma)f(\tilde\gamma),\qquad
 z(\tilde\gamma)=P(\tilde\gamma)^2f(\tilde\gamma),
\]
where $P(\tilde\gamma)$ is a unit. Conditionally on $a$ and $c$, $f(\tilde\gamma)$ is
uniform on $J$, because the constant coefficient of $f$ has a free uniform $J$-part.
Require all these evaluations to be nonzero, at a cost of at most $q/M$. At the other
residue points $y,z$ have nonzero reductions. Consequently $y$ has no root among the
chosen representatives, while $z$ has exactly one, namely $\tilde\beta$.

One further condition identifies the start from $y$. Write $Y=\bar y$. For every simple
root $\gamma$ of $Y$ test the equality in $J$
\[
 y(\tilde\gamma)=-y_n\,\frac{\bigl(Y/(x-\gamma)\bigr)(\gamma)}{\gpplc(Y)} .
\]
It holds at $\beta$. At any other simple root of $Y$ the right side is fixed after
conditioning on $a,c$, and the left side is uniform on $J$, as just proved. Discard the
starts with a second successful test, at a cost of at most $q/M$. All these costs use the
union bound on the original starts and require no independence between the tests.

The resulting triples are disjoint. Degree distinguishes starts from both other
positions, and evaluation at the representatives distinguishes middles from ends. A
middle identifies $\beta$ by its unique successful test; then $a=\bar y/(x-\beta)$ and
$c=-y_n/\gpplc(a)$ recover $r$. An end identifies $\beta$ by its unique root and gives
$a=\bar z/(x-\beta)^2$; dividing in $A[x]$ by the monic $P$ yields
$(z/P)(\tilde\beta)=2ca(\beta)$ and hence $c$. Finally $r$ is a non-zero-divisor, since
$r(P-c(P^s+1))=P^2$ is monic, so each inversion recovers at most one start. The ordinary
triple belongs to the family \eqref{gpp:eq:pattern} for each $L\ge2$, and removing one
member of each disjoint triple gives the lower bound for $1-\Delta_L(A)$.
\end{proof}

\subsection{Finite coefficient rings: square-zero algebras and natural density}\label{gpp:sec:squarezero}
{\raggedright\sources{grh_method_transfers_research; grh_method_transfers_graph_towers_and_finite_rings}}

For square-zero extensions of a field the squares of nonunits can be described exactly.
This gives sharper bounds for ordinary triples, exact extrema for a related pair
condition, and a strict separation between natural and upper density.

Let $k=\Fq$, let $V$ be a $k$-vector space of dimension $b-1$, and give $R=k\oplus V$ the
multiplication $(u,v)(u',v')=(uu',uv'+u'v)$. Put $Q=\abs{R}=q^b$. With a fixed basis of
$V$ write $f\in R[x]$ as a tuple $(f_0,f_1,\ldots,f_{b-1})$ with entries in $k[x]$. In
this subsection a set is \emph{admissible} if it is contained in $R[x]\setminus\set{0}$
and contains no triple $f,rf,r^2f$ with $r\in R[x]$ a nonunit; retaining zero would
change no density. An element $r=(a,c)$ is a unit exactly when $a\in k^\times$, as
$(a,c)^{-1}=(a^{-1},-a^{-2}c)$ in that case. Let $\Delta(R)$ be the supremum of the upper
densities of admissible sets and $D(R)$ the supremum of the natural densities of
admissible sets whose ball densities converge.

\begin{theorem}[Characteristic two; \stU]\label{gpp:thm:char2}
Suppose $q$ is even.
\begin{enumerate}[label=\textup{(\roman*)}]
\item The set $S$ of nonzero tuples whose monic componentwise gcd is squarefree is
admissible; it consists exactly of the nonzero elements divisible by no square of a
nonunit. For every $n\ge2$,
\[
 \abs{S\cap B_n}=Q^{n+1}-qQ^{n-1}+q-1,
\]
so $S$ has natural density $1-q/Q^2$.
\item One has
\begin{equation}\label{gpp:eq:localring}
 1-\frac1{Q^2}\le\Delta(R)\le1-\frac{Q-1}{Q^3-1}=1-\frac1{Q^2+Q+1}.
\end{equation}
For fixed $q$ this determines the leading constant: $Q^2(1-\Delta(R))\to1$ as
$b\to\infty$. The upper bound holds in every characteristic and for $b=1$.
\item Among sets containing no pair $f,a^2f$ with $a\in k[x]$ nonconstant, the largest
upper density is exactly $1-Q^{-2}$ and the largest natural density is exactly
$1-q/Q^2$; the two extrema differ by $(q-1)/Q^2$.
\end{enumerate}
All these assertions hold more generally for a finite local $k$-algebra $R=k\oplus V$
with nilpotent maximal ideal $V$, provided every $v\in V$ satisfies $v^2=0$; distinct
elements of $V$ may then have nonzero product.
\end{theorem}

\begin{proof}
In characteristic two $r=a+c$ with $a\in k[x]$, $c\in V[x]$ has $r^2=a^2$, since
Frobenius acts coefficientwise on $c$ and all its coefficients square to zero. In the
more general algebra the unit criterion follows by reduction modulo the nilpotent ideal
$V[x]$ and a finite geometric series. If $a=0$ the last term of a triple is zero.
Otherwise choose a monic irreducible $P\mid a$; every component of the last term is
divisible by $P^2$. Conversely every nonzero tuple with a common square factor $P^2$ is a
multiple of the square of the scalar nonunit $P$. This proves the description of $S$ and
its admissibility.

For a nonzero tuple the indicator that its gcd is squarefree is
$\sum_{d^2\mid\gcd(f_i)}\mu(d)$ over monic $d\in k[x]$. In $B_n$ there are
$Q^{n+1-2j}-1$ nonzero multiples of a fixed $d^2$ of degree $2j\le n$. Unique
factorization gives $\sum_{d\text{ monic}}\mu(d)u^{\deg d}=1-qu$, the reciprocal of
$\sum_{j\ge0}q^ju^j$. Thus only $j=0,1$ contribute, which yields the count and the
natural density in (i).

Next consider sets containing no pair $f,a^2f$ with $a\in k[x]$ nonconstant; every such
set is admissible. In a single ball $B_n$, $n\ge1$, the two layers of highest degrees
$n-1,n$ avoid all such pairs and have $Q^{n+1}-Q^{n-1}$ elements. For the reverse bound
map each nonzero $f\in B_n$ to $x^{2\floor{(n-\deg f)/2}}f$ in those two layers. Two
distinct preimages of one point differ by multiplication by a positive even power of
$x$, so the map is injective on any pair-avoiding set.

To attain the limiting proportion in one infinite set, let $h(f)$ be the largest monic
polynomial whose square divides every component of $f$; at each prime its exponent is
half the common valuation rounded down. Thus $m(f)=\deg f-2\deg h(f)$ is invariant under
multiplication by any nonzero scalar polynomial square. Choose $N_0=-1$ and $N_j=2^{j+1}$
for $j\ge1$. Assign each $f\ne0$ to the unique $j$ with $N_{j-1}<m(f)\le N_j$, and
include it exactly when $\deg f\in\set{N_j-1,N_j}$. The members of a forbidden pair would
have the same assigned $j$ and degrees differing by at least two, which is impossible. In
the top two layers of $B_{N_j}$ a missing element has $m(f)\le N_{j-1}$, hence
$\deg h(f)\ge L_j=\ceil{(N_j-1-N_{j-1})/2}$. For a monic $d$ of degree $s\le N_j/2$ the
fraction of $B_{N_j}$ divisible componentwise by $d^2$ is $Q^{-2s}$, and a union bound
shows that the missed fraction is at most $\sum_{s\ge L_j}(q/Q^2)^s$, which tends to
zero. This proves the lower bound in \eqref{gpp:eq:localring} and the value $1-Q^{-2}$
in (iii). The construction also applies for $b=1$.

For the natural-density extremum in (iii), each nonzero $f$ factors uniquely as $uh^2$
with $h=h(f)$ monic and $u\in S$. For a pair-avoiding set $E$ let $I_u$ be the set of
monic $h$ with $uh^2\in E$. If $1\in I_u$ there are no other elements. Otherwise
\[
 \sum_{h\in I_u}Q^{-2\deg h}\le\sum_{s\ge1}(q/Q^2)^s\le1,
\]
since $q/Q^2\le1/2$. In either case the mass is at most one. Multiplying by
$Q^{-\deg u}$ and summing over the base tuples up to degree $n$ gives
\[
 \sum_{f\in E\cap B_n}Q^{-\deg f}\le\sum_{u\in S\cap B_n}Q^{-\deg u}.
\]
The number of all tuples of exact degree $t$ is $(Q-1)Q^t$. Whenever a set has natural
ball density $\delta$, its exact-degree proportions tend to $\delta$ as well, by
subtracting consecutive ball counts. Dividing the last inequality by $(n+1)(Q-1)$ and
using the convergence of arithmetic means bounds the natural density of $E$ by
$1-q/Q^2$, and $S$ attains it.

For the upper bound in \eqref{gpp:eq:localring} take a scalar monic linear $P$. Every
nonzero tuple has an exact common $P$-valuation, the minimum of the component valuations
with $v_P(0)=+\infty$. Partition each chain $h,Ph,P^2h,\ldots$ with common valuation of
$h$ equal to zero into triples starting at exponents $0,3,6,\ldots$. Multiplication by
$P$ is injective and increases the maximum component degree by one. Hence the triples
are disjoint, and every triple lying in $B_n$ forces at least one deletion. Their number
is exactly
\[
 T_n=\sum_{0\le3j\le n-2}\bigl(Q^{n-1-3j}-Q^{n-2-3j}\bigr),
\]
since their first terms have degree at most $n-2$ and valuation $3j$. Dividing by
$Q^{n+1}$ and summing the limiting geometric series gives
$(1-Q^{-1})Q^{-2}/(1-Q^{-3})$. This argument works in every characteristic, including
$b=1$.
\end{proof}

\begin{theorem}[A strict gap for triple avoidance; \stU]\label{gpp:thm:gaprings}
Suppose $q$ is even. Then
\[
 1-\frac q{Q^2}\le D(R)\le1-\frac q{Q^2}(1-Q^{-1})^q .
\]
Thus $Q^2(1-D(R))\to q$ as $b\to\infty$ with $q$ fixed, and $D(R)<\Delta(R)$ for every
$b\ge2$.
\end{theorem}

\begin{proof}
The lower bound is \cref{gpp:thm:char2}(i). For the upper bound let $C$ consist of the
nonzero tuples whose gcd has no linear factor. By reduction modulo the product of the $q$
monic linears, the natural density of $C$ is $(1-Q^{-1})^q$. For each $g\in C$ form the
$q$ triples $g,Pg,P^2g$, one for each monic linear $P$. Two different centres give
disjoint collections, because the common linear valuations of a member determine the
multiplying polynomial $1$, $P$ or $P^2$ uniquely; within one collection only the centre
$g$ is shared. Give a tuple $f$ of degree $d$ the weight $Q^{-d}$. Omitting $g$ costs
$Q^{-\deg g}$. If $g$ is present, at least one of $Pg,P^2g$ is missing for each $P$,
which costs at least $qQ^{-2-\deg g}$. In either case the omitted weight is at least
$qQ^{-2-\deg g}$. Sum over the centres in $B_{n-2}$, so that every tuple counted is in
$B_n$, divide by $(n+1)(Q-1)$, and pass to the limit using the exact-degree proportions
of $C$ and of the omitted set as in the proof of \cref{gpp:thm:char2}. This gives the
upper bound. For $b\ge2$ one has $q(1-Q^{-1})^q>q(1-q/Q)\ge q-1\ge1$ by the strict
Bernoulli inequality, so the natural-density upper bound is strictly below the
upper-density construction $1-Q^{-2}$.
\end{proof}

\begin{theorem}[Odd characteristic; \stU]\label{gpp:thm:oddnatural}
If $q$ is odd, the nonzero elements of $R[x]$ divisible by no square of a nonunit form an
admissible set of natural density
\begin{equation}\label{gpp:eq:localringodd}
 \prod_{P\text{ monic irreducible}}
 \bigl(1-z_P^{b+1}+z_P^{b+2}-z_P^{2b+1}\bigr)\ge1-q^{-b},\qquad z_P=q^{-\deg P}.
\end{equation}
\end{theorem}

\begin{proof}
Every forbidden triple ends at a multiple of a nonunit square. Discard the tuples with
$f_0=0$, a set of proportion $q^{-(n+1)}$. A tuple with $f_0\ne0$ is a multiple of the
square of a nonunit if and only if some prime $P$ satisfies one of the two mutually
exclusive conditions
\[
\begin{array}{ll}
 v_P(f_0)=2, &P\mid f_i\quad(1\le i<b),\\
 v_P(f_0)\ge3, &P^2\mid f_i\quad(1\le i<b).
\end{array}
\]
For necessity write $f=(a,c)^2(h_0,h)=(a^2h_0,a^2h+2ach_0)$ and take $P\mid a$. All
nilpotent coordinates are multiples of $P$. If $v_P(a^2h_0)\ge3$, then $P^2\mid a$ or
$P\mid h_0$, and all are multiples of $P^2$. For sufficiency, in the second case divide
by the pure square $P^2$. In the first case set $h_0=f_0/P^2$. Since both $2$ and $h_0$
are invertible modulo $P$, choose $c_i$ solving $2c_ih_0\equiv f_i/P\pmod P$. Then
$h_i=(f_i-2Pc_ih_0)/P^2$ is a polynomial for each $i$ and $f=(P,c)^2(h_0,h)$.

For fixed $P$, uniform residues of independent polynomial coordinates modulo $P^3$ give
the total local obstruction probability
$(z_P^2-z_P^3)z_P^{b-1}+z_P^3z_P^{2(b-1)}$. The Chinese remainder theorem gives
independence for any finite set of primes once $n$ is sufficiently large. For nonzero
$f_0$ an obstruction with $\deg P=d$ requires $2d\le n$ and has probability at most
$q^{-(b+1)d}$, because it needs $P^2\mid f_0$ and $P\mid f_i$ for all $i>0$. Thus the
omitted large primes have total probability at most $\sum_{d>D}q^{-bd}$, which tends to
zero uniformly in $n$. This gives the product. Since $2b+1\ge b+2$ for $b\ge1$, every
factor is at least $1-z_P^{b+1}$, and the product of these lower bounds is $1-q^{-b}$ by
unique factorization.
\end{proof}

\begin{example}[Dual numbers]\label[example]{gpp:ex:dual}
For $R=\F_2[\epsilon]/(\epsilon^2)$, so $q=2$ and $Q=4$, the squarefree construction has
natural density $7/8$; every admissible set with a natural density has density at most
$119/128$; the block construction has upper density $15/16$; and
$15/16\le\Delta(R)\le20/21$. More generally any even $q$ and $b\ge2$ give a lower bound
at least $1-q^{-3}$, already for natural density, which exceeds the upper bound
$1-1/(q^2+q+1)$ valid for the residue field itself (the case $b=1$ of
\eqref{gpp:eq:localring}; it is $6/7$ for $\F_2[x]$), because $q^3>q^2+q+1$. By
\eqref{gpp:eq:localringodd} the same comparison holds for arbitrary $q$ in square-zero
extensions of dimension $b\ge3$. In the opposite direction, for the same $q$ and $b$ the
reduced algebra $\Fq^b$ has the same cardinality $Q$, yet $\Delta(\Fq^b)\to0$ by
\cref{gpp:thm:local}: its coordinate idempotents permit the independent ratios that
realize the ternary word cube.
\end{example}

\subsection{Periodic sets with a prescribed spectrum of ratios}\label{gpp:sec:syndetic}
{\raggedright\sources{dirichlet-grh-research}}

Beiglb\"ock, Bergelson, Hindman and Strauss asked whether every set of positive integers
with bounded gaps contains arbitrarily long geometric progressions
\cite{Beiglbock2006Multiplicative}; even three terms with an integral common ratio
$r\ge2$ are at issue \cite{Nathanson2020Geometric}. Consider the finite restriction in
which a prime modulus $p$ is fixed and ratios outside a prescribed residue class are to
be excluded. The weighted congruence of \cref{gpp:lem:digits} is exactly the relation
between the three discrete logarithms of a triple $a,ar^i,ar^j$, and the lemma gives the
exact set of permissible ratios for many exponent pairs at once.

For a nonempty periodic set $W\subset\N$ write $G(W)$ for the largest difference between
successive elements and $d(W)$ for the natural density. For a positive rational $r=u/v$
in lowest terms, its residue $\bar r=u/v\in\F_p^\times$ is defined when $p\nmid uv$.

\begin{theorem}[Exact ratio spectrum of a periodic set; \stU]\label{gpp:thm:syndetic}
Let $p$ run through the primes, and let $L=L(p)\ge2$ be integers with
$\log L=o(\log p)$. Put $E=\sqrt{\log p\,\log(2L)}$. For all sufficiently large $p$
there is a $p$-periodic set $W\subset\N$, all of whose members are coprime to $p$, with
\[
 d(W)\ge\exp(-CE),\qquad G(W)\le\sqrt p\,\exp(CE)
\]
for an absolute constant $C$, such that for every $1\le i<j\le L$ and every positive
rational $r=u/v$ in lowest terms, some integer $a$ satisfies $a,ar^i,ar^j\in W$ if and
only if
\[
 p\nmid uv\qquad\text{and}\qquad\bar r^{\gcd(i,j)}=1 .
\]
Whenever a ratio occurs, infinitely many $a$ realize it.
\end{theorem}

In particular the ordinary triple $a,ar,ar^2$ occurs exactly for $r\equiv1\pmod p$ among
rational $p$-units, and the least integral ratio greater than one is $p+1$. A single
occupied residue class would have the same ratio spectrum; the theorem gives
$p^{1-o(1)}$ occupied classes and an additive gap of at most $p^{1/2+o(1)}$,
simultaneously for all the specified shapes.

\begin{lemma}[\stU]\label[lemma]{gpp:lem:gap}
Every $m$-element set $B\subset\F_p$, $m\ge1$, has a nonzero multiplicative dilate with
maximum circular gap at most $2(\floor{p/\sqrt m}+1)$.
\end{lemma}

\begin{proof}
Set $h=\floor{p/\sqrt m}+1$. If $2h\ge p$ the assertion is automatic. Otherwise take
independent uniform $x\in\F_p$ and $t\in\F_p^\times$, and let $Z$ count the points of
$tB+x$ in $I=\set{0,\ldots,h-1}$. Each image is uniform, and any two distinct images are
uniformly distributed as an ordered pair of distinct points. Consequently
$\mu=\gppE Z=mh/p$ and $\gppVar(Z)\le\mu$, since the hit indicators have nonpositive
pairwise covariance. Chebyshev's inequality gives $\gppPr(Z=0)\le1/\mu$. Averaging in
$t$, some dilate has at most $p/\mu=p^2/(mh)<h$ empty intervals of length $h$ among the
$p$ cyclic starting positions. A run of $2h-1$ vacant points would contain $h$ such
intervals.
\end{proof}

\begin{proof}[Proof of \cref{gpp:thm:syndetic}]
Apply \cref{gpp:lem:digits} with $o=p-1$ and
\[
 n=\floor*{\sqrt{\frac{\log o}{\log(2L)}}},\qquad H=\floor*{\frac{o^{1/n}}{L}} .
\]
Since $\log L=o(\log p)$, both $n$ and $o^{1/n}/L$ tend to infinity. Eventually
$o^{1/n}/(2L)\le H\le o^{1/n}$ and $(LH)^n\le o$, and the size $m=\abs{S}$ satisfies
\[
 m\ge\frac{o}{n(2L)^no^{2/n}}\ge o\exp\bigl(-O(\sqrt{\log o\,\log(2L)})\bigr),
\]
because the logarithm of the first denominator is $\log n+n\log(2L)+2(\log o)/n$.

Choose a generator $g$ of $\F_p^\times$ and write $B=\set{g^e:e\in S}$. Any
multiplicative dilate $tB$ may be used; let $W$ be the inverse image of $tB$ in $\N$. If
$a,ar^i,ar^j\in W$, their $p$-adic valuations vanish and $r$ is a $p$-unit. Write their
residues as $tg^{e_0},tg^{e_i},tg^{e_j}$. Then
\[
 je_i\equiv(j-i)e_0+ie_j\pmod{p-1}.
\]
The lemma forces the three exponents to agree. Thus $\bar r^i=\bar r^j=1$, which is
equivalent to $\bar r^{\gcd(i,j)}=1$. Conversely, under these conditions choose $a$
divisible by $v^j$ in any class of $tB$; the Chinese remainder theorem gives infinitely
many choices, and all three integers have the same residue modulo $p$. Finally take for
$t$ the dilate of \cref{gpp:lem:gap}. Then $d(W)=m/p$ and $G(W)\ll p/\sqrt m$.
\end{proof}

\begin{proposition}[Obstruction to a fixed gap bound; \stU]\label[proposition]{gpp:prop:syndeticobstruction}
For any nonempty $p$-periodic set, all positive integral ratios $r\equiv1\pmod p$ occur
in progressions of arbitrary length. Let $W$ be a $p$-periodic set in which all integral
ratios of ordinary triples are $1\pmod p$. Then $W$ omits the residue zero, the discrete
logarithms of its occupied residues, represented in $[0,p-2]$, contain no nonconstant
three-term arithmetic progression, and
\[
 d(W)\le\exp\bigl(-c(\log p)^{1/9}\bigr),\qquad G(W)\ge\exp\bigl(c(\log p)^{1/9}\bigr)
\]
for large $p$ and an absolute $c>0$. If moreover $W$ has the simultaneous spectra of
\cref{gpp:thm:syndetic} for some $L<p-1$ and occupies $m$ residue classes, then
$m(L+1)\le2(p-1)$ and $G(W)\ge p(L+1)/(2(p-1))$.
\end{proposition}

\begin{proof}
If $W$ contained a multiple of $p$, the multiples of $p$ in $W$ would realize every
integral ratio. Exponents $e,e+d,e+2d$ of occupied residues give a triple of residues
with ratio $g^d\ne1$, and a positive integer representative of this ratio gives a triple
in $W$. The density bound then follows from the bound of Bloom and Sisask for sets
without three-term arithmetic progressions \cite{Bloom2023Improvement}, and the maximum
gap is at least the mean gap. For the last assertion, three occupied exponents among
$L+1$ consecutive cyclic exponents would give a forbidden triple with ratio $g$ and
$j\le L<p-1$; counting occupied exponents over all such intervals gives
$m(L+1)\le2(p-1)$.
\end{proof}

Thus increasing the prime period cannot yield a fixed gap bound by this method. McNew
already treated geometric progressions in modular unit groups through arithmetic
progressions in a finite abelian group \cite[Section~2]{McNew2015Sets}; the deduction
here combines the simultaneous shape control of \cref{gpp:lem:digits}, an additive gap
estimate, and exact realization of every admissible rational ratio.

\begin{example}\label[example]{gpp:ex:syndetic}
\gppsloppy
For $p=65\,537$, $g=3$, $L=4$, $n=4$ and $H=4$, a largest digit sphere has squared radius
$14$ and $24$ vectors. Searching all dilations gives $t=18\,054$ and largest circular gap
$4\,555$. Among the six pairs $1\le i<j\le4$, the shape $(0,2,4)$ admits the residues $1$
and $-1$; each of the other five admits only $1$. The program
\artifact{dirichlet-grh-research/progressions/verify_syndetic.py} checks all $393\,222$
combinations of exponent pair and residue ratio, including zero, and
\artifact{dirichlet-grh-research/progressions/syndetic-example.txt} records the occupied
residues.
\end{example}

\subsection{Sparse places in the source construction}\label{gpp:sec:sparse}
{\raggedright\sources{exponent_seven_and_grh_transfers; mcnew_counterexample_and_grh_transfers}}

The choice of infinite places in the source proof poses a finite extremal problem on
residue classes. Fix a positive multiple $q_0$ of $18480$ and an integer $s\ge1$, and put
\[
 z=3^{\floor{\log_3s}+4},\quad q_1=q_0\cdot3z\cdot2^s,\quad v=v_2(q_0),\quad
 C=\set{x\in(\Z/q_1\Z)^\times:x\equiv-1\ (\mathrm{mod}\ q_0)} .
\]
The source identifies the place labels with $(\Z/q_1\Z)^\times$, and the desired labels
lie in $C$. For a prime $p$ let $m_p=q_1/p^{v_p(q_1)}$. A set $T\subseteq C$ is
\emph{sparse at $p$} if there are no $x,y,t\in T$, integers $0<i<j<z$ and signs
$\epsilon,\delta\in\set{1,-1}$ with
\[
 y/x\equiv\epsilon p^i\pmod{m_p},\qquad t/x\equiv\delta p^j\pmod{m_p}.
\]
The three entries may coincide and the signs are independent. This is the local
condition in \code{New.Parameters.IsSparse}; imposing it at all primes $p\le L$ is the
sparse-set requirement of the source. The source theorem
\code{Parameters.exists_large_isSparse} uses $(2zH)^n\le2^s$ and $2z(L+1)<2^s$ to obtain
at least $H^n/(nH^2)$ sparse places for the primes $p\le L$; it fixes one odd residue and
uses a large set of exponents in the $2$-primary coordinate.

\subsubsection*{Which primes can be admitted}

Let $\delta_p$ denote the order of the class of $p$ in $(\Z/m_p\Z)^\times/\set{\pm1}$,
that is, the least positive $t$ with $p^t\equiv\pm1\pmod{m_p}$.

\begin{theorem}[Exact obstruction for sparse places; \stU]\label{gpp:thm:sparseexact}
Let $H\ge2$ and $n\ge1$ satisfy $\bigl((z-1)(H-1)+1\bigr)^n\le2^s$. There is a single set
$T\subseteq C$ with
\[
 \abs{T}\ge\frac{H^n}{nH^2}
\]
such that for every prime $p$ the following are equivalent: $T$ is sparse at $p$; some
nonempty subset of $C$ is sparse at $p$; and $2\delta_p\ge z$.
\end{theorem}

\begin{proof}
A prime with $2\delta_p<z$ obstructs every nonempty set: use $x=y=t$, $i=\delta_p$ and
$j=2\delta_p$. We construct a set for which all obstructions arise in this way.

Let $S\subset\set{0,\ldots,2^s-1}$ be the set of \cref{gpp:lem:digits} for $o=2^s$ and
coefficient budget $L=z-1$, so that every relation
$je_2\equiv(j-i)e_1+ie_3\pmod{2^s}$ with $0<i<j<z$ and $e_1,e_2,e_3\in S$ forces
$e_1=e_2=e_3$. Put $g=1+3zq_0$ and $T=\set{-g^e\bmod q_1:e\in S}$. The integer $g$ is
prime to $q_1$ and congruent to $1$ modulo $q_0$. As $v_2(g-1)=v\ge4$, one has
$v_2(g^a-1)=v+v_2(a)$ for $a>0$: an odd power preserves the valuation, and each
successive squaring adds one. Hence $g$ has order $2^s$ modulo $N=2^{s+v}$, and the
labels are distinct.

For odd $p$, the divisor $N$ of $m_p$ retains both congruences of any forbidden triple.
On $(\Z/N\Z)^\times$ send a unit $u$ to its unique signed representative
$\rho(u)\in\set{u,-u}$ that is $1\pmod4$. This is a group homomorphism, it kills both
signs, and it fixes $g$. Writing the exponents of $x,y,t$ as $e_1,e_2,e_3$, the
congruences give
\[
 g^{e_2-e_1}=\rho(p)^i,\qquad g^{e_3-e_1}=\rho(p)^j\pmod N .
\]
Raise the first relation to the power $j$ and the second to the power $i$. The order of
$g$ gives $je_2\equiv(j-i)e_1+ie_3\pmod{2^s}$, so $e_1=e_2=e_3$. Returning to the full
modulus $m_p$, both $i$ and $j$ are multiples of $\delta_p$, so $2\delta_p\le j<z$.

For $p=2$, every ratio of two elements of $T$ is $1$ modulo $3z$, a divisor of $m_2$. A
forbidden triple would imply $4^i\equiv1\pmod{3z}$. The order of $4$ modulo $3z$ is $z$,
since $v_3(4^a-1)=1+v_3(a)$ for $a>0$; this contradicts $i<z$. Hence $T$ is sparse at
$2$, and by the first paragraph $2\delta_2\ge z$.
\end{proof}

\begin{corollary}[The largest prime cutoff; \stU]\label[corollary]{gpp:cor:cut}
The largest integer $L$ for which some nonempty subset of $C$ is sparse at every prime
$p\le L$ is
\[
 \min\set{p\text{ prime}:2\delta_p<z}-1,
\]
and the set $T$ of \cref{gpp:thm:sparseexact} realizes it. In particular
$z(L+1)\le2^{s+v+1}$ is sufficient.
\end{corollary}

\begin{proof}
A prime $p\equiv1\pmod{q_1}$ has $\delta_p=1$, and such primes exist by Dirichlet's
theorem; this gives a first obstruction. For odd $p$ choose $a\in\set{p,-p}$ with
$a\equiv1\pmod4$. Its order modulo $N$ is $2^{\max(0,s+v-v_2(a-1))}$ by the valuation
formula above, and this order divides $\delta_p$. Since $2^{v_2(a-1)}\le p+1$, we have
$\delta_p\ge2^{s+v}/(p+1)$. For $p=2$ the proof of \cref{gpp:thm:sparseexact} applies.
\end{proof}

\begin{example}[A $120$-place instance]\label[example]{gpp:ex:120}
Take $q_0=18480$, $s=44$ and $z=2187$, so that
\[
 q_1=2^{48}\cdot3^9\cdot5\cdot7\cdot11=2\,133\,004\,707\,139\,399\,188\,480 .
\]
With $n=3$ and $H=12$ the digit base is $24047$ and $24047^3=13\,905\,375\,151\,823<2^{44}$.
Exactly $120$ triples $(a,b,c)\in\set{0,\ldots,11}^3$ satisfy
$(2a-11)^2+(2b-11)^2+(2c-11)^2=131$: the unordered triples of odd squares are
$(121,9,1)$, $(81,49,1)$ and $(81,25,25)$, and signs and coordinate permutations give
$48+48+24$ labels $-g^{a+24047b+24047^2c}\bmod q_1$.

The sufficient cutoff of the source is $4\,021\,990\,406$, the elementary bound of
\cref{gpp:cor:cut} gives $257\,407\,386\,108$, and the exact largest cutoff is
$18\,216\,708\,648\,992\,766$; the first obstructing prime is
$p_*=18\,216\,708\,648\,992\,767$. Indeed, for an odd obstructing prime one must have
$\delta_p\le1093$, and its projective order modulo $2^{48}$ divides $\delta_p$. This
dyadic order is a power of two at most $1024$, so $p\equiv\pm1\pmod{2^{38}}$. Of the
$132\,543$ integers at least $2$ in these two classes through $p_*$, exactly $55\,083$
are prime to $q_1$, and the primes dividing $q_1$ other than $2$ fail the two-class test.
For each retained integer $u$, compute its dyadic order and test its multiples $k$
through $1093$ for $u^k\equiv\pm1\pmod{q_1}$. Exactly seven integers pass.
\begin{center}
\begin{tabular}{rrl}
\toprule
Integer & Least exponent & Certificate\\
\midrule
$1\,020\,346\,790\,576\,129$ &1080& divisible by $157$\\
$2\,185\,829\,116\,018\,687$ &320& divisible by $101999$\\
$3\,870\,280\,929\,771\,521$ &972& divisible by $1663$\\
$6\,412\,351\,813\,189\,631$ &864& divisible by $17$\\
$9\,399\,724\,905\,857\,023$ &768& divisible by $23$\\
$17\,099\,604\,835\,172\,353$ &324& divisible by $13$\\
$18\,216\,708\,648\,992\,767$ &864& Lucas certificate, base $3$\\
\bottomrule
\end{tabular}
\end{center}
For the last integer, $p_*-1=2\cdot3^6\cdot11\cdot74413\cdot15264089$ with prime factors
verified by trial division, $3^{p_*-1}\equiv1\pmod{p_*}$, and
$\gcd(3^{(p_*-1)/\ell}-1,p_*)=1$ for each prime factor $\ell$; this proves primality. At
this prime $i=864$ and $j=1728$ obstruct even a singleton. The computation is in
\artifact{exponent_seven_and_grh_transfers/other-results/verify_sparse.py}, with output
\artifact{exponent_seven_and_grh_transfers/other-results/sparse_output.txt}.
\end{example}

\subsubsection*{How many places can be used}

\Cref{gpp:thm:sparseexact} keeps one odd residue. Two of the prime constraints admit an
exact joint solution, which bounds the capacity from above and enlarges the
construction.

For $p\in\set{2,3}$ write $b_p=q_0/p^{v_p(q_0)}$ and $N_p=m_p/b_p$, so that $N_2=3z$ and
$N_3=2^s$. Let $d_p$ be the least positive integer with $p^{d_p}\equiv\pm1\pmod{b_p}$,
and let $g_p=\epsilon_pp^{d_p}$, where $\epsilon_p\in\set{-1,1}$ is the unique sign
giving $g_p\equiv1\pmod{b_p}$. Define
\[
 H_p=\ord_{m_p}(g_p),\qquad w_p=\ceil{z/d_p},\qquad
 A_p=\frac{N_p}{H_p}\min\Bigl(H_p,\floor*{\frac{2H_p}{w_p}}\Bigr).
\]
Both $d_p$ depend only on $q_0$.

\begin{theorem}[The exact two-prime capacity; \stU]\label{gpp:thm:sparsetwo}
The largest size of a set $T\subseteq C$ that is sparse at both $2$ and $3$ is $A_2A_3$.
Thus $A_2A_3$ is the exact optimum for the prime cutoffs $L=3,4$ and an upper bound for
every $L\ge3$; in particular
\[
 \frac{\abs{T}}{\abs{C}}\le\frac{4d_2d_3}{z^2}\qquad(L\ge3).
\]
\end{theorem}

\begin{proof}
The possible images in $(\Z/m_p\Z)^\times$ of elements of $C$ are precisely the residues
$a\equiv-1\pmod{b_p}$. They number $N_p$, since every prime divisor of $m_p$ already
divides $b_p$. If $\pm p^i$ takes one of these residues to another, then $d_p\mid i$ by
the order of the class of $p$ in $(\Z/b_p\Z)^\times/\set{\pm1}$. Since $b_p>2$ the sign
is unique, and the step is multiplication by $g_p^{i/d_p}$. Multiplication by $g_p$
partitions the possible images into $N_p/H_p$ cycles of length $H_p$.

Write $H=H_p$ and $w=w_p$. On one such cycle mark the residues in the image of $T$, and
read their indicator as a periodic binary sequence on $\Z$. Three marked positions in a
window of $w$ consecutive indices give a forbidden triple, by subtracting the first index
and multiplying both differences by $d_p$. Conversely the three indices of any forbidden
triple fit in such a window, because $\floor{(z-1)/d_p}=w-1$. If a window goes around the
cycle more than once, all its occurrences count, as required for repeated places. For
$K$ marks the sum of the occupancies of the $H$ windows is $wK$, so $K\le H$ and
$K\le\floor{2H/w}$. Both bounds are attained together: with $K=\min(H,\floor{2H/w})$
mark the indices $\ell$ with
\[
 \floor*{\frac{(\ell+1)K}{H}}-\floor*{\frac{\ell K}{H}}=1 .
\]
There are $K$ such indices per period, and the sum over every window is at most
$\ceil{wK/H}\le2$ by telescoping. Hence the largest allowable projected support has
$A_p$ elements.

Write $q_0=2^v3^ab$ with $\gcd(b,6)=1$. The component modulo $b$ is fixed to $-1$ in
every element of $C$. The two remaining freely varying components are a lift of
$-1\bmod2^v$ and a lift of $-1\bmod3^a$. The condition at $2$ uses only the second of
these, and the condition at $3$ only the first. Any two permissible projected supports
can therefore be combined by taking their full Cartesian product under the Chinese
remainder theorem, and every eligible $T$ is contained in the product of its own
projections. This proves exactness. Finally $A_p/N_p\le2/w_p\le2d_p/z$ and
$\abs{C}=N_2N_3$.
\end{proof}

\begin{example}[The base modulus $18480$]\label[example]{gpp:ex:18480}
Here $b_2=1155$ and $b_3=6160$. The orders of $2$ modulo $3,5,7,11$ are $2,4,3,10$, and
the powers of $2$ cannot reach $-1$ modulo $7$. The orders of $3$ modulo $16,5,7,11$ are
$4,4,6,5$, and the powers of $3$ cannot reach $-1$ modulo $16$. Consequently
$d_2=d_3=60$ and both signs are positive. The identities $v_3(2^{60}-1)=2$ and
$v_2(3^{60}-1)=4$ give $H_2=z$ and $H_3=2^s$: replacing the exponent $60$ by $60k$ adds
$v_3(k)$ to the first valuation and $v_2(k)$ to the second, since
$3^{60}\equiv1\pmod4$, and the other prime-power components already have order one. With
$w=\ceil{z/60}$ the capacity is therefore
\[
 3\min\bigl(z,\floor{2z/w}\bigr)\,\min\bigl(2^s,\floor{2^{s+1}/w}\bigr).
\]
For $s\ge81$ this equals $357\floor{2^{s+1}/\ceil{z/60}}$. Indeed $z=6561$ gives $w=110$
and $\floor{2z/w}=119$; the next possible value is $z=19683$, and for it and all larger
$z$ one has $120z/(z+59)>119$ while $2z/w<120$ because $60\nmid z$. The $3z$ possible
odd residues have thus shrunk to exactly $357$ before any prime $p\ge5$ is considered.
\end{example}

\begin{proposition}[An all-prime construction; \stU]\label[proposition]{gpp:prop:sparseall}
Let $n,H$ be positive integers and $L\ge2$ an integer with
\[
 \bigl((z-1)H\bigr)^n\le2^s,\qquad(L+1)(z-1)<2^{v+s}.
\]
There exists $T\subseteq C$, sparse at every prime $p\le L$, with
\[
 \abs{T}\ge A_2\,\frac{H^n}{nH^2}.
\]
\end{proposition}

\begin{proof}
Let $S\subset\set{0,\ldots,2^s-1}$ be as in the proof of \cref{gpp:thm:sparseexact}. Take
as the $2$-primary component $-(1+2^v)^e\bmod2^{v+s}$ for $e\in S$, and independently an
optimal support of size $A_2$ in the odd modulus $m_2$, as in \cref{gpp:thm:sparsetwo}.
The first components are distinct, since $1+2^v$ has order $2^s$. The condition at $2$
holds by construction. Let $p\le L$ be odd and let $t_p\in\set{p,-p}$ be $1\bmod4$.
Project a proposed forbidden triple to $2^{v+s}$, a divisor of $m_p$. The two ratios are
$1\bmod4$, so their signs are forced and their values are $t_p^i$ and $t_p^j$. Raising
the first to the $j$th power and the second to the $i$th gives the weighted relation on
exponents, so all three exponents are equal and $t_p^i\equiv1\pmod{2^{v+s}}$. Yet
\[
 2^{v_2(t_p^i-1)}=2^{v_2(t_p-1)}2^{v_2(i)}\le(p+1)i<2^{v+s},
\]
a contradiction.
\end{proof}

Compared with the source theorem, for the same $n$ and $H$ the lower bound is multiplied
by $A_2$ and both size conditions are relaxed; the coefficient budget drops from $2z$ to
$z-1$ because the sign is determined modulo $4$ before the powers are compared.
\Cref{gpp:thm:sparseexact} admits a still larger range of primes, namely exactly the
primes that any nonempty set can admit, but keeps one odd residue. For $q_0=18480$,
$s=40$, $L=19$, $n=2$ and $H=40$ the construction of \cref{gpp:prop:sparseall} takes
eight exponents on one digit sphere and $354$ projected residues at $2$, giving $2832$
places.

\begin{proposition}[Sharp order at a fixed prime cutoff; \stU]\label[proposition]{gpp:prop:sparseorder}
Let $M(q_0,s,L)$ be the largest size of a subset of $C$ that is sparse at every prime
$p\le L$. For fixed $q_0$ and fixed $L\ge3$, as $s\to\infty$,
\[
 M(q_0,s,L)=\Theta_{q_0,L}(2^s/z),\qquad
 \frac{M(q_0,s,L)}{\abs{C}}=\Theta_{q_0,L}(z^{-2}).
\]
More precisely, if $r$ is the number of odd primes up to $L$ and
$(L+1)(z-1)<2^{v+s}$, then
\[
 \frac{A_2\,2^s}{1+2r(z-1)}\le M(q_0,s,L)\le A_2A_3 .
\]
\end{proposition}

\begin{proof}
The upper bound is \cref{gpp:thm:sparsetwo}. For the lower bound form a graph on the
$2^s$ possible $2$-primary components, joining $x$ to $y$ if $y/x$ or $x/y$ is a signed
power $p^i$ for some odd prime $p\le L$ and $0<i<z$. The valuation bound in the proof of
\cref{gpp:prop:sparseall} rules out loops once $(L+1)(z-1)<2^{v+s}$, which holds for all
large $s$ since $z\le81s$. For each pair $p,i$ the residue modulo $4$ forces the sign, so
the degree is at most $2r(z-1)$. Greedily choosing a vertex and discarding its
neighbours yields an independent set of size at least $2^s/(1+2r(z-1))$. This excludes
even the first nontrivial step of each forbidden triple at an odd prime. Use its full
product with an optimal support at the prime $2$. The latter has $A_2\ge1$: a singleton
odd residue could return under a signed $2^i$ only if $4^i\equiv1\pmod{3z}$, whereas
$4=1+3$ has order $z$ there.
\end{proof}

{\gppsloppy The program
\artifact{mcnew_counterexample_and_grh_transfers/checks/sparse_places.py} evaluates the
two orders, enumerates all binary cycle supports for $108$ small pairs of size and
window, and tests explicit attaining supports against the signed-power predicate.\par}

\subsection{Coefficient collisions and higher-dimensional anticodes}\label{gpp:sec:lattice}
{\raggedright\sources{grh_method_transfers_2026-11-26}}

In \cref{gpp:lem:digits} a prescribed family of coefficient relations between exponents
is excluded. One may change which relations must be excluded and ask how small an
encoding group suffices. Allow every equality between two sums of $h$ labels, with
repetitions, as a potential collision. A label set free of such collisions between
distinct multisets is a classical $B_h$ set. A coefficient difference
$z=(z_0,\ldots,z_n)$ has zero total and $\tfrac12\sum_i\abs{z_i}\le h$ after common
labels cancel; conversely each such difference comes from two $h$-term sums by padding
both with the same label. Thus the coefficient geometry is that of the root lattice
\[
 A_n=\set{z\in\Z^{n+1}:\textstyle\sum_iz_i=0},\qquad
 d(z,w)=\tfrac12\textstyle\sum_i\abs{z_i-w_i},
\]
a classical correspondence \cite{Kovacevic2017Improved}.

Write $a_n(D)$ for the largest cardinality of a subset of $A_n$ of diameter at most $D$.
For $x\in\Z^n$ put $\pi(x)=\sum_i\max(x_i,0)$ and $\nu(x)=\pi(-x)$. The isometry
$x\mapsto(x,-\sum x_i)$ identifies $d$ with the norm $\max(\pi,\nu)$ of the difference.
Put $P=\ceil{D/2}$ and $Q=\floor{D/2}$. The balanced difference of two discrete
simplices $S_n(P,Q)=\set{x:\pi(x)\le P,\ \nu(x)\le Q}$ has diameter at most $P+Q$, since
$\pi(x-y)\le\pi(x)+\nu(y)$ and likewise for $\nu$. Choosing the $j$ strictly positive
coordinates and using stars and bars on each sign gives its cardinality
\[
 M_n(D)=\sum_{j=0}^{\min(n,P)}\binom nj\binom Pj\binom{Q+n-j}{n-j}.
\]
Kova\v{c}evi\'c conjectures that $a_n(D)=M_n(D)$ for all $n,D$ and proves this for $n=3$
\cite[\S7]{Kovacevic2026Isodiametric}.

\begin{theorem}[Exact anticodes in ranks four and five; computer-assisted; \stU]\label{gpp:thm:anticodes}
The values of $a_4(D)$ for $2\le D\le10$ and of $a_5(D)$ for $2\le D\le6$ are
\begin{center}
\begin{tabular}{@{}r*{9}{r}@{}}\toprule
$D$ & 2 & 3 & 4 & 5 & 6 & 7 & 8 & 9 & 10\\\midrule
$a_4(D)$ & 21 & 55 & 131 & 255 & 471 & 780 & 1251 & 1875 & 2751\\
$a_5(D)$ & 31 & 96 & 271 & 606 & 1281 & & & &\\\bottomrule
\end{tabular}
\end{center}
In each case $a_n(D)=M_n(D)$, and every maximizing family is a translate of $S_n(P,Q)$ or
of its negative.
\end{theorem}

\begin{proof}
Let $F$ be a set of diameter at most $D$ in root coordinates and put $m=n+1$. Subtract
each coordinate's minimum on $F$. All coordinates now lie in $[0,D]$ and all coordinate
sums equal some $k$; the reflection $x_i\mapsto D-x_i$ allows
$0\le k\le\floor{mD/2}$. Each possible ambient layer $V_k$ is finite. For $x,y$ in it,
\[
 d(x,y)=\max_{J\subseteq\set{1,\ldots,m}}\sum_{i\in J}(x_i-y_i),
\]
since the sum of all positive coordinate differences is half their total absolute value.
Accordingly $F$ has diameter at most $D$ precisely when each subset sum has range at most
$D$ on $F$. Singletons already satisfy this, and a cut and its complement have equal
range.

For each remaining cut let $\alpha,\beta$ be its minimum and maximum on the whole layer.
If $\beta-\alpha>D$, try each integer interval $[\sigma,\sigma+D]$ for
$\alpha\le\sigma\le\beta-D$; any feasible range is contained in one of these intervals,
by clipping its lower endpoint to $\beta-D$ when necessary. If $\beta-\alpha\le D$, the
entire layer satisfies the cut. Intersect with a chosen interval and recurse, and stop
when the intersection has fewer than $M_n(D)$ points. Every feasible family fits some
branch at every depth, so the absence of a larger intersection on all full branches
proves the upper bound, and at equality the family is the entire intersection. For every
leaf of size at least $M_n(D)$ the program checks its size and compares the full set,
after subtracting coordinate minima, with the two normalized candidate shapes. This
proves the classification of equality. \gppsloppy The program
\artifact{grh_method_transfers_2026-11-26/checks/lattice/anticodes.py} performs these
operations with exact integer bit sets; its outputs
\artifact{grh_method_transfers_2026-11-26/checks/lattice/A4.txt} and
\artifact{grh_method_transfers_2026-11-26/checks/lattice/A5.txt} record the node counts.
\end{proof}

The two entries for $D=2$ also follow from classical Singer difference sets for the
prime powers $4$ and $5$ through the code--anticode bound
\cite{Kovacevic2017Improved,Kovacevic2026Isodiametric}. The theorem for rank three does
not cover either row for $D\ge3$.

\begin{corollary}[Multisets with large pairwise intersections; \stU]\label[corollary]{gpp:cor:multisets}
Let $(n,D)$ be one of the pairs of \cref{gpp:thm:anticodes} and $m=n+1$. A family of
$K$-multisets on $m$ symbols in which every two members share at least $K-D$ elements,
counted with multiplicity, has size at most $M_n(D)$, with equality possible exactly
when $K\ge mQ+P-Q$. For example, a family of $K$-multisets on five symbols in which
every pair shares at least $K-3$ elements has size at most $55$, with equality exactly
when $K\ge6$.
\end{corollary}

\begin{proof}
For two count vectors of sum $K$ the intersection cardinality is $K-d(x,y)$, so the
upper bound holds on every layer. For the lower bound take $x\in S_n(P,Q)$ and add the
final coordinate $P-Q-\sum_ix_i$. All $m$ coordinates are at least $-Q$ and their sum is
$P-Q$. Each coordinate attains $-Q$, so the normalized shape is obtained by adding $Q$
everywhere; it has coordinate sum $mQ+P-Q$. The negative shape interchanges $P$ and $Q$.
Nonnegative translations give all larger sums. The classification of equality therefore
gives the threshold $K\ge mQ+P-Q$, since $P\ge Q$ and $m\ge2$.
\end{proof}

\begin{proposition}[Approximation in every fixed rank; \stU]\label[proposition]{gpp:prop:anticodeasymptotic}
For every fixed $n$, $a_n(D)=M_n(D)+O_n(D^{n-1})$.
\end{proposition}

\begin{proof}
Let $\mathcal B=\set{x\in\R^n:\pi(x),\nu(x)\le1}$, with $\pi,\nu$ extended to real
vectors. Attaching unit cubes in the $\Z^n$ coordinates to a family of size $a_n(D)$
gives a union $U$ of that volume with $U-U\subseteq(D+n)\mathcal B$. The
Brunn--Minkowski inequality gives
$a_n(D)\le2^{-n}(D+n)^n\operatorname{vol}(\mathcal B)$. The formula for $M_n(D)$ has
leading coefficient $2^{-n}\binom{2n}{n}/n!$, which equals
$2^{-n}\operatorname{vol}(\mathcal B)$: sum the orthant volumes
$\binom nj/(j!(n-j)!)$ over $j$.
\end{proof}

\begin{theorem}[The least order of a three-sum encoding with five labels; computer-assisted; \stU]\label{gpp:thm:sidon}
The least order of a finite abelian group containing a $B_3$ set of cardinality five is
$64$, and the least order among cyclic groups is $65$. Witnesses are
\[
 \begin{array}{ll}
 (\Z/8\Z)^2: & (0,0),(0,1),(1,0),(4,6),(7,6),\\
 \Z/65\Z: & 0,1,4,15,24 ,
 \end{array}
\]
each giving $\binom73=35$ distinct unordered sums of three labels. In particular a
full-rank sublattice of $A_4$ of minimum distance at least four has index at least $64$,
and there is no diameter-perfect lattice code in $A_4$ at $D=3$.
\end{theorem}

\begin{proof}
Translate a candidate $\set{b_0,\ldots,b_4}$ in a group $G$ so that $b_0=0$. The
homomorphism $A_4\to G$, $z\mapsto\sum_{i=0}^4z_ib_i$, has no nonzero kernel element of
distance at most $3$ from the origin, by cancellation and padding. It is therefore
injective on the $55$ points of $S_4(2,1)$ in root coordinates, and $\abs{G}\ge55$. For
the orders $55$ through $63$, the structure theorem gives the cyclic group at every
order and just four additional groups, with invariant factors $(2,28)$, $(2,2,14)$,
$(2,30)$ and $(3,21)$.

The program \artifact{grh_method_transfers_2026-11-26/checks/lattice/sidon.py} orders
each group lexicographically and enumerates all increasing four-label extensions of
$\set{0}$. It retains a partial set $T$ exactly when its unordered triple sums are
distinct. When a label $x$ is adjoined, the new sums are the elements of
$(x+2T)\cup(2x+T)\cup\set{3x}$, where $2T$ is the set of pairwise sums; there must be
precisely $\abs{2T}+\abs{T}+1$ new values, all outside $3T$. This test counts all new
multisets correctly, because the partial $B_3$ property implies that pairwise sums are
distinct, by adding one fixed old label. Removing a label preserves the property, so
every valid set is reached. All thirteen groups of order $55$ to $63$ fail, and so does
the cyclic group of order $64$; the output is
\artifact{grh_method_transfers_2026-11-26/checks/lattice/sidon.txt}. Both displayed
witnesses pass an additional direct check over the $35$ multisets.

For the last assertion, the five coordinate labels in the quotient of $A_4$ by such a
sublattice, after translating a common coordinate vector to zero, have the $B_3$
property. A diameter-perfect lattice code at $D=3$ would have index $a_4(3)=55$.
\end{proof}

\subsection{Open questions}\label{gpp:sec:open}
{\raggedright\sources{grh-methods-polynomial-density; grh-progression-densities-and-cubic-prime-race; mcnew-density-and-dirichlet-grh; grh_method_transfers; grh_method_transfers_graph_towers_and_finite_rings; grh_method_transfers_research; dirichlet-grh-research; exponent_seven_and_grh_transfers; mcnew_counterexample_and_grh_transfers; grh_method_transfers_2026-11-26}}

\begin{question}[The size of $\beta_q-m_q$]\label{gpp:q:beta}
Determine $\beta_q$. For $q=2$ the results above give
$2^{-2^{476}}<\beta_2-m_2<5.46\cdot10^{-7}$; computing the gain to the first decimal
place at which $\beta_2$ differs from $m_2$ is a specific numerical problem, which a
smaller degree substitution or a sharper optimized smooth part in
\cref{gpp:lem:replacement} could advance, as could the solution of the finite extremal
problems of \cref{gpp:thm:supremum}. For large $q$, \cref{gpp:thm:largeq} bounds the
difference by $q^{-q+O(1)}$ up to a polynomial factor; is there a lower bound on a
comparable scale? This would require a much more economical construction than that of
\cref{gpp:thm:counterexample}.
\end{question}

\begin{question}[Rational-function ratios]\label{gpp:q:rational}
Determine the supremum of the upper densities of rational-ratio progression-free subsets
of $\Fq[x]$. \Cref{gpp:thm:rational} computes the optimum when all coordinates but one
are subject to the greedy restriction. Optimizing two free valuation coordinates
requires avoiding collinear midpoint triples also in mixed directions.
\end{question}

\begin{question}[Radicals of higher nilpotence]\label{gpp:q:cube}
For finite local rings with $\gppm^3=0$ and $\gppm^2\ne0$, which invariant controls the
loss of density? The elimination \eqref{gpp:eq:eliminate} then leaves a term in the
second radical layer; whether its rank or its multiplication law enters the criterion is
the first case beyond \cref{gpp:thm:local}. More generally, classify the growing products
of arbitrary finite local rings with positive limiting density, and determine the
extremal densities for general finite coefficient rings.
\end{question}

\begin{question}[Products of a fixed field]\label{gpp:q:decay}
For fixed $q$, determine the rate of decay of $\Delta(\Fq^b)$ as $b\to\infty$. The
density Hales--Jewett reduction gives vanishing and \cref{gpp:prop:behrendrings} gives a
lower bound $\exp(-C\sqrt{\log(2+b)})$. Can polynomial ratios beyond those with entries
$1,x$ substantially sharpen the upper bound?
\end{question}

\begin{question}[Exact extrema for small rings]\label{gpp:q:dual}
Determine $\Delta(R)$ and $D(R)$ for the dual numbers
$R=\F_2[\epsilon]/(\epsilon^2)$, where $15/16\le\Delta(R)\le20/21$ and
$7/8\le D(R)\le119/128$. Does allowing both endpoints of a square-ratio pair increase
the achievable upper density beyond $15/16$? Determine the exact extrema for finite
products of fields, for instance for $\F_2\times\F_2$.
\end{question}

\begin{question}[Syndetic sets]\label{gpp:q:syndetic}
Does every set of positive integers with bounded gaps contain a three-term geometric
progression with an integral ratio $r\ge2$? By \cref{gpp:prop:syndeticobstruction},
increasing the prime period in \cref{gpp:thm:syndetic} cannot yield a fixed gap bound, so
the question for an arbitrary syndetic set and unrestricted integral ratios is not
reached by that construction.
\end{question}

\begin{question}[Capacity of the sparse-place condition]\label{gpp:q:sparse}
Determine $M(q_0,s,L)$ exactly for $L\ge5$. How does the maximal density of sparse
places decay when the prime cutoff grows with the parameter $s$, as in the source proof?
The proof of \cref{gpp:prop:sparseorder} tracks the dependence on the number of primes,
and the exact two-prime extremizers of \cref{gpp:thm:sparsetwo} are specific candidates
on which to test the remaining congruences.
\end{question}

\begin{question}[Anticodes in $A_n$]\label{gpp:q:anticodes}
Is $a_n(D)=M_n(D)$ for all $n\ge4$ and all $D$, with the translates of the balanced
simplex difference and of its negative as the only extremal families?
\Cref{gpp:prop:anticodeasymptotic} gives only the leading term.
\end{question}
